
\documentclass[11pt]{article}
\usepackage[a4paper,margin=24mm,headheight=15pt]{geometry}
\usepackage[T1]{fontenc}
\usepackage[utf8]{inputenc}
\usepackage{lmodern,amsmath,amssymb,amsthm,mathtools}
\usepackage{booktabs,longtable,array}
\usepackage{xcolor,hyperref,microtype,fancyhdr,enumitem}
\hypersetup{colorlinks=true,linkcolor=blue!45!black,urlcolor=blue!45!black,citecolor=blue!45!black,
pdftitle={Shared factorial errors in a Diophantine equation for Con(ZFC)},
pdfauthor={OpenAI ChatGPT, prepared for Vladimir Reshetnikov}}
\pagestyle{fancy}\fancyhf{}
\fancyhead[L]{\small Shared factorial errors and Diophantine consistency}
\fancyhead[R]{\small 25 October 2026}
\fancyfoot[C]{\thepage}
\setlength{\parskip}{5pt}\setlength{\parindent}{0pt}
\setlist{itemsep=3pt,topsep=4pt}
\allowdisplaybreaks[4]
\newtheorem{theorem}{Theorem}[section]
\newtheorem{lemma}[theorem]{Lemma}
\newtheorem{proposition}[theorem]{Proposition}
\newtheorem{corollary}[theorem]{Corollary}
\theoremstyle{definition}\newtheorem{definition}[theorem]{Definition}
\newcommand{\Z}{\mathbb Z}\newcommand{\N}{\mathbb N}
\newcommand{\Con}{\operatorname{Con}}\newcommand{\ZFC}{\mathrm{ZFC}}
\newcommand{\ZF}{\mathrm{ZF}}\newcommand{\Ccalc}{\mathsf C222}
\newcommand{\Gform}{\mathcal G}
\newcommand{\abs}[1]{\left|#1\right|}
\newcommand{\MainVars}{240}
\newcommand{\MainLast}{239}
\newcommand{\MainRows}{177}
\newcommand{\MainRowLast}{176}
\newcommand{\MainTerms}{21482}
\newcommand{\MainHeight}{6144}
\newcommand{\MainDir}{equation240}
\title{\textbf{Shared factorial errors in a Diophantine equation for
\(\Con(\ZFC)\)}\\[5pt]\large A direct quartic with \MainVars{} unrestricted
integer unknowns}
\author{Prepared for Vladimir Reshetnikov and the ProveIt project\\
\small Mathematical development and computational assistance: OpenAI ChatGPT}
\date{25 October 2026}
\begin{document}\maketitle
\begin{abstract}
We give a completely specified degree-four polynomial in \MainVars{}
unrestricted integer unknowns whose lack of integer zeros is provably
equivalent in ZFC to \(\Con(\ZFC)\). The final equation has \MainRows{}
quadratic residuals, \MainTerms{} collected monomials, and coefficient
height \MainHeight{}. Every residual is printed in this report.

The main mathematical change shares one norm certificate between eleven
modular digit remainders and a signed factorial approximation error.
A nearest-integer factorial lemma recovers the intended factorial after
binomial extraction, making two former positivity certificates unnecessary.
Centered digits also allow an initially signed odd radix, whose positivity
is recovered later. The original beta-spacing formula is retained.
These changes give a 244-variable direct core; exact integral elimination
and reversible residual-row operations produce the stated final equation.

The route uses the numerical proof calculus C222 and its identified ZF/ZFC
refutation calibration, without instantiating K215 or Jones's universal
polynomial. We distinguish ordinary unbounded proofs from finite algebra
checks. No global simplicity minimum, consistency proof, or newly completed
end-to-end proof-assistant formalization is claimed.
\end{abstract}
\newpage\tableofcontents\newpage
\section{An explicit equation and the scope of the result}

The principal deliverable is the following single polynomial equation over
unrestricted integers:
\begin{equation}\label{eq:main-definition}
 P(x_0,\ldots,x_{\MainLast})
   :=\sum_{i=0}^{\MainRowLast}Q_i(x_0,\ldots,x_{\MainLast})^2=0.
\end{equation}
The appendix gives every coefficient of every quadratic polynomial \(Q_i\).
The file \texttt{\MainDir/\allowbreak equation.json} is the same finite definition in a
sparse, machine-readable format; the accompanying collected coefficient
file gives \(P\) directly. There are no unspecified universal polynomials,
variable-exponent operations, or machine simulations in this definition.
The powers appearing in the proof below describe the meanings of ordinary
polynomial registers.

\begin{theorem}\label{thm:main}
For the fixed integer polynomial in \eqref{eq:main-definition}, ZFC proves
\[
 \bigl(\forall \mathbf X\in\mathbb Z^{\MainVars}\quad P(\mathbf X)\ne0\bigr)
       \quad\Longleftrightarrow\quad \operatorname{Con}(\mathrm{ZFC}).
\]
More precisely, ZFC proves that an integer zero exists if and only if a
finite refutation exists in the fixed numerical presentation of ZFC used
by the retained ProveIt calibration.
\end{theorem}

The equation has \(\MainVars\) variables, \(\MainRows\) squared quadratic
residuals, total degree four, \(\MainTerms\) nonzero collected monomials,
and coefficient height \(\MainHeight\). The height is the largest absolute
coefficient after all like monomials have been collected. All variables
range over \(\mathbb Z\), including coordinates whose eventual arithmetic
interpretations are nonnegative.

The comparison is with the incoming common-remainder construction at the
pinned repository snapshot, rather than with an older, larger ancestor.
The exact endpoints supplied here make the tradeoffs visible:
\begin{center}
\begin{tabular}{lrrrr}
\hline
Equation & Variables & Residuals & Monomials & Height\\
\hline
Common-remainder parent &252&182&5340&128446\\
New arithmetic core &244&181&5346&128446\\
Primary dimension endpoint &\MainVars&\MainRows&\MainTerms&\MainHeight\\
Balanced companion &241&178&5351&15107\\
Lower-height companion &241&178&9544&4512\\
\hline
\end{tabular}
\end{center}
All five entries have degree four. The primary endpoint gives the smallest
dimension among the new supplied forms, at a substantial cost in support:
its collected polynomial has about four times as many terms as the
252-variable parent. The balanced companion keeps almost the parent's
support while removing eleven variables and reducing coefficient height.
These are distinct exact equations with the same existence theorem.

This is an improvement of the direct numerical proof-calculus branch. It
does not establish a global minimum. Other branches of the repository
already contain quartics in far fewer variables, including a 37-integer
form, with different coefficient and encoding costs \cite{overview}.
There is no established scalar measure called \emph{simplest} that orders
dimension, degree, coefficient size, support, and proof complexity.
Here the declared objective is an explicit quartic with small coefficients,
a direct refutation interpretation, and a complete reproducible description.

The new arithmetic step is to use the existing common norm to certify the
factorial error as well as all eleven digit remainders. This removes the
two initial factorial sign certificates. A signed-radix proof then permits
the radix certificate to be only an affine oddness equation. These changes
are compatible with the inherited deterministic Pell comparison center.
The common normalized remainders, the center construction, and the
underlying calculus are credited to their predecessors
\cite{common,symmetric,centered,repo}; the proof below identifies precisely
which additional implications are needed.

\section{The retained arithmetic interface}

Write \(t=16+d^2\), where \(d\) is an unrestricted integer. The principal
arithmetic quantities are
\[
 F,\quad B,\quad V,\quad x,\quad y,\quad N,\quad q,\quad
 M,\qquad H=1+\sum_{\nu=1}^{9}a_\nu.
\]
The \(a_\nu\) are the nine actual CRT representatives. In the stored source,
several are a coordinate plus one; it is their actual values, rather than
their shifted coordinates, that will become nonnegative.

The retained exact-power requests have the intended meanings
\begin{equation}\label{eq:power-meanings}
 V=(F^2+1)^t,\qquad N=B^x,\qquad q=(BH)^x,\qquad M=q^x,
\end{equation}
together with the powers used in the binomial numerators and denominators.
There are eighteen requests, sharing fourteen independent principal Pell
pairs. At this point these displays are requested meanings, not assumptions.
Their interpretation is recovered in a specified order below.

The beta spacing is left literally unchanged:
\begin{equation}\label{eq:beta}
 y=V(V-1).
\end{equation}
The endpoint equations include
\begin{equation}\label{eq:endpoints}
 x=F+1+(y+1)u_0,\qquad
 x=1+\bigl(1+(t+1)y\bigr)u_t.
\end{equation}
Initially \(u_0,u_t,x,F,B,q\) may have either sign.

The eleven normalized remainder registers fall into four groups:
\begin{center}
\begin{tabular}{llr}
\hline
Denominator & Meaning & Number\\
\hline
\(q^t\) & Factorial digit and encoded index &2\\
\(q^{t+1}\) & Four syntax representatives and derived label &5\\
\(q^{y+1}\) & Two pending-multiset representatives &2\\
\(q^x\) & Two beta quotients &2\\
\hline
\end{tabular}
\end{center}
The existing ratio equations are
\begin{equation}\label{eq:ratios}
 q^t v_t=M,\qquad q^{t+1}v_{t+1}=M,\qquad
 q^{y+1}v_y=M.
\end{equation}
Each displayed power in these equations is a retained output register.
The quotient for the last group is the constant one.

For clarity, put
\[
 G(a,b,c)=a^2+b^2+c(c+1).
\]
This form is nonnegative at every integer triple and represents every
nonnegative integer. The latter is the familiar consequence of the
three-square theorem applied to \(4n+1\): exactly one square is odd, and
the two even squares can be divided by four. This is the same guard form
already used in the parent \cite{common}.

The altered arithmetic core replaces the three separate initial scalar
sign certificates by
\begin{align}
 q&=2z+1,                                                    \label{eq:odd}\\
 Bv&=M,                                                     \label{eq:factor-ratio}\\
 MF&=vV+S,                                                  \label{eq:factor-error}\\
 M^2&=1+4\left(S^2+\sum_{i=1}^{11}R_i^2\right)+G(a,b,c).
                                                               \label{eq:common-norm}
\end{align}
The new integers \(v,S\) normalize the factorial error at the same
denominator \(M\) as the other remainders. No inequality, division,
factorial, or variable power has been added to the literal polynomial
system. Equations \eqref{eq:odd}--\eqref{eq:common-norm} are polynomial
rows of degree at most two in their actual registers.

\section{Pell orientation before arithmetic signs}

For a signed integer center \(A\), define the Pell polynomials by
\[
 X_0(A)=1,\quad Y_0(A)=0,\qquad
 X_{k+1}=AX_k+(A^2-1)Y_k,\quad Y_{k+1}=X_k+AY_k.
\]
For negative indices use \(X_{-k}=X_k\), \(Y_{-k}=-Y_k\).
The retained frame bounds every actual base, requested index, and proposed
power output in absolute value by an integer \(W\), with
\begin{equation}\label{eq:frame}
 W\ge258,\qquad |A|>W^W,\qquad |A-1|>4W^3.
\end{equation}
These are consequences of square-based frame equations. Their proof does
not assume signs for \(F,B,q,x\), or the power outputs.

The deterministic comparison center is retained from the supplied
predecessor proof \cite{symmetric}. To distinguish its auxiliary pair
from the arithmetic quantity \(H\), write it as
\[
 H_{\rm p}^2-(A^2-1)J_{\rm p}^2=1,\qquad
 J_{\rm p}=C^2(2e+1),
\]
where \((U,C)\) is the principal anchor pair with requested index \(h\).
The squared cap says
\[
 C^2=h^2+1+\sum_{\text{other pairs}}Y_i^2+G(a',b',c').
\]
It gives \(C\ne0\), \(|h|<|C|\), and strict domination of all other
principal ordinates. The comparison center is defined by
\[
 f=(1-A)H_{\rm p},\qquad D=A+H_{\rm p}f
    =1-(A-1)(A^2-1)J_{\rm p}^2.
\]
Thus \(|D|>1\), \(D\equiv A\pmod{H_{\rm p}}\), and \(D\equiv1\pmod C\).
The comparison ordinate residue is \(R=h+Ck\).

Here is the part of its proof that matters for compatibility. Classify
the anchor and common Pell ordinates at \(|A|\) by positive indices
\(i,l\). For \(i\ge2\), the squared divisibility condition implies
\(|C|\mid l\) and \(l>i\). The abscissa comparison gives the comparison
index \(j\equiv\pm i\pmod{2l}\); the new ordinate residue gives
\(j\equiv\pm h\pmod{|C|}\). Consequently
\(i\equiv\pm h\pmod{|C|}\). The growth bound
\[
 |C|=Y_i(|A|)\ge i+2|A|-2
\]
and \(|h|<W\) imply \(i+|h|<|C|\), so \(i=|h|<W\).
All other classified indices are smaller by the cap.

The exceptional case \(i=1\) has \(|C|=1\). The cap then forces \(h=0\)
and every other ordinate to vanish. Reducing the retained orientation
pool modulo \(A-1\) gives a nonnegative small error sum. Its anchor
ordinate contributes one, so it is strictly positive, while
\eqref{eq:frame} makes it strictly smaller than \(|A-1|\).
Its asserted divisibility is impossible. This exceptional branch
therefore supplies no zero.

All fourteen principal pairs retain both-coordinate orientation:
the squared abscissa errors and squared ordinate-index errors remain
in the pool. After classification their small residues have errors at
most two and strictly less than \(2W\), respectively. The resulting
nonnegative residue sum is less than
\(4\cdot14(W^2+1)<4W^3<|A-1|\). Divisibility forces every error to vanish.
Thus each principal pair is exactly
\((X_k(A),Y_k(A))\) at its signed requested index.

This use of all fourteen pairs is deliberate. A separate predecessor
reduction keeps fewer abscissa tests, but its hypotheses cannot simply
be assumed in the signed bootstrap below. The present construction
does not make that deletion. The complete retained center and frame
proofs are included in the proof closure.

\section{A power lemma valid at signed bases}

\begin{lemma}\label{lem:signed-power}
Assume \eqref{eq:frame}. Let \(|n|,|k|,|m|<W\), with \(|n|\ge2\), and set
\[
 T=2An-n^2-1.
\]
If the pair is canonically oriented at \(k\) and its power row asserts
\[
 X_k(A)+(n-A)Y_k(A)\equiv m\pmod T,
\]
then \(k\ge0\) and \(m=n^k\).
For \(n=\epsilon\in\{1,-1\}\), the same row instead gives
\(m=\epsilon^k\) for every integer \(k\), including negative \(k\).
\end{lemma}

\begin{proof}
For \(|n|\ge2\),
\[
 |T|>4W^W-(W-1)^2-1>W^W+W.
\]
Modulo \(T\), one has
\((n-A)^2=A^2-1\) and \(n(2A-n)=1\).
The Pell identities therefore give \(n^k\) at nonnegative indices
and its multiplicative inverse at negative indices. If \(k\ge0\),
the congruence has
\[
 |m-n^k|<W+W^W<|T|,
\]
so it is equality. If \(k=-j<0\), multiplication by \(n^j\) gives
\(mn^j\equiv1\pmod T\). But
\[
 |mn^j-1|\le(W-1)^W+1<|T|,
\]
forcing the impossible integer equality \(mn^j=1\).

For a unit base, \(T=2A\epsilon-2\) and \(2A\equiv2\epsilon\pmod T\).
The sequence
\(f_k=X_k(A)+(\epsilon-A)Y_k(A)\) has initial values
\(f_0=1,f_1=\epsilon\) and satisfies
\(f_{k+1}=2Af_k-f_{k-1}\). This recurrence runs backwards as well
as forwards. Hence \(f_k\equiv\epsilon^k\pmod T\) for every signed
integer \(k\). Finally \(|T|>2W^W-2>W\), whereas
\(|m-\epsilon^k|\le W\), giving the stated equality.
\end{proof}

The base \(-1\) case is essential: it cannot be replaced by a lemma
that presupposes a nonnegative exponent. The literal request for
\(M=q^x\) has the row
\[
 -M+X_x(A)-AY_x(A)+qY_x(A)-T_q\eta=0.
\]
Both \(M\) and \(q\), and the actual base \(BH\) of the self-radix
request, are included in the frame. Thus the lemma applies to the
unshifted arithmetic values claimed here.

\section{One norm also certifies the factorial error}

The shared-radius argument works before the denominator signs are known.
Suppose nonzero integers \(L,v_i,M\) satisfy \(Lv_i=M\) and
\[
 v_i U=M Q+R_i.
\]
Then \(v_i\mid R_i\); setting \(r_i=R_i/v_i\) gives an integer and
recovers \(U=LQ+r_i\). Equation \eqref{eq:common-norm} implies
\[
 |R_i|<|M|/2,\qquad |r_i|<|L|/2.
\]
If a prescribed signed remainder \(r_i^0\) is congruent to \(U\)
modulo \(L\) and has \(|r_i^0|<|L|/2\), uniqueness follows:
the multiple \(r_i-r_i^0\) of \(L\) has absolute value less than \(|L|\).
Positivity of either \(L\) or \(r_i^0\) is unnecessary.

The same mechanism yields the new factorial estimate.
First, \eqref{eq:common-norm} gives \(M\ne0\).
Equations \eqref{eq:factor-ratio}--\eqref{eq:factor-error} then imply
\[
 B\ne0,\quad v\ne0,\quad S=v(BF-V).
\]
With \(b=|B|\) and \(s=BF-V\), this gives
\begin{equation}\label{eq:near}
 s\in\mathbb Z,\qquad |s|<b/2.
\end{equation}
The canonical pair at the already positive index \(t\) gives
\(V=(F^2+1)^t\ge1\), using the unit case if \(F=0\).

\begin{lemma}\label{lem:early-bound}
Under these conditions, \(b\ge2\), \(b>|F|\), and \(V<b^2\).
If \(F=0\), in fact \(b\ge3\). If \(F\ne0\), \(B\) and \(F\)
have the same sign.
\end{lemma}

\begin{proof}
When \(F=0\), \(V=1\) and \(|BF-V|=1<b/2\), so \(b>2\).
Suppose \(f=|F|\ge1\). Opposite signs would give
\(|BF-V|=bf+V>b/2\), a contradiction. Therefore \(BF=bf\).
If \(b\le f\), \eqref{eq:near} would give
\[
 V<b(f+1/2)\le f^2+f/2,
\]
contradicting \(V\ge(f^2+1)^2>f^2+f/2\).
Thus \(b\ge f+1\ge2\), and
\[
 V<b(f+1/2)\le b(b-1/2)<b^2.
\]
\end{proof}

This estimate supplies the coefficient margin formerly obtained from
two separate nonnegativity graphs. It has not yet identified \(B\)
as a binomial coefficient or \(F\) as a factorial.

\section{Recovering the radix and the first exact digit}

\subsection{Excluding the two unit radices}

Equation \eqref{eq:odd} makes \(q\) odd, hence nonzero. If \(q=1\)
or \(q=-1\), Lemma~\ref{lem:signed-power} applied to the request
\(M=q^x\) gives \(|M|=1\), with no assumption on \(x\).
But \(Bv=M\) and \(|B|\ge2\) make this impossible. Consequently
\[
 Q:=|q|\ge3,\qquad |q+1|\ge2.
\]
Every actual representative \(a_\nu\) is a requested exponent at
base \(q+1\). The same lemma now gives \(a_\nu\ge0\), and therefore
\(H\ge1\). The derived label \(h\) is also an exponent at that base,
so \(h\ge0\); the anchor conclusion excludes \(h=0\).
The request at base \(q\) gives \(x\ge0\) and \(M=q^x\).

The retained ratio \(q^{t+1}v_{t+1}=M\) is now an equality of
nonzero integer powers. Its absolute values imply \(x\ge t+1\).
Since \(|BH|\ge b\ge2\), the self-radix request gives
\begin{equation}\label{eq:signed-radix-size}
 q=(BH)^x,\qquad Q\ge b^x\ge b^{t+1}.
\end{equation}
Neither the radix nor \(B\) has been assumed positive in this argument.
The beta value \eqref{eq:beta} is nonnegative, but its size is not
needed to reach \eqref{eq:signed-radix-size}.

\subsection{A binomial coefficient at a possibly negative radix}

Put \(n=F^2+t\) and \(C_j=\binom{n}{j}\). The exact numerator request
at base \(q+1\) gives \(U=(q+1)^n\). Dividing its normalized digit
row by its nonzero ratio gives
\begin{equation}\label{eq:factor-digit}
 U=q^t(B+q\eta)+r,\qquad |r|<Q^t/2.
\end{equation}
For \(F\ne0\), every coefficient through position \(t\) satisfies
\begin{equation}\label{eq:coeff}
 C_j\le(F^2+1)^t=V<b^2\qquad(0\le j\le t).
\end{equation}
For completeness, the general inequality
\(\binom{m+t}{j}\le(m+1)^t\), for \(m\ge1\) and \(j\le t\),
follows by induction on \(t\). Pascal's identity handles
\(j\le t\) in the next row using \(2\le m+1\); the new edge is
bounded by
\(\binom{m+t+1}{t+1}=\prod_{i=1}^{t+1}(1+m/i)\le(m+1)^{t+1}\).
If \(F=0\), use the separate bound \(C_j\le2^t\).
Using \(V=1\) to bound those coefficients would be incorrect.

For either sign of \(q\), the binomial theorem gives the signed
low polynomial
\[
 r^0=\sum_{j=0}^{t-1}C_jq^j,\qquad
 |r^0|<\frac{C_{\max}}{Q-1}Q^t.
\]
When \(F\ne0\), \eqref{eq:signed-radix-size} gives
\[
 \frac{C_{\max}}{Q-1}
 <\frac{2}{b^{t-1}}<\frac1{32}.
\]
When \(F=0\), \(b\ge3\), and
\[
 \frac{C_{\max}}{Q-1}
 \le\frac{2^t}{3^{t+1}-1}
 <(2/3)^{t+1}<\frac1{32}
 \qquad(t\ge16).
\]
Thus \(|r^0|<Q^t/32\). It is congruent to \(U\) modulo
the possibly negative integer \(q^t\). The signed remainder
uniqueness argument gives \(r=r^0\).

Subtract the low polynomial from \eqref{eq:factor-digit},
divide by \(q^t\), and reduce modulo \(q\). The result is
\[
 B\equiv C_t\pmod q.
\]
In the nonzero-\(F\) case, \(|B-C_t|\le b+C_t<b+b^2<Q\).
In the zero-\(F\) case, \(C_t=1\) and \(b+1<Q\).
Therefore in both cases
\begin{equation}\label{eq:B-exact}
 B=\binom{F^2+t}{t}>0.
\end{equation}
The case \(F=0\) would give \(B=1\), contradicting \(b\ge3\).
Hence \(F\ne0\); the sign conclusion of
Lemma~\ref{lem:early-bound} now gives \(F>0\).
Finally \(q=(BH)^x\) with \(B>0,H\ge1,x\ge t+1\) gives \(q>0\).
The negative-radix branch has served its purpose and has been excluded
by the arithmetic itself. All remaining exact-power requests, including
\(N=B^x\), now have their usual positive-base meanings.

\section{The nearest-integer factorial fixed point}

The new bound is symmetric around zero, so the old ceiling lemma alone
does not justify the construction. The required strengthened statement
has a short direct proof.

\begin{lemma}\label{lem:nearest-factorial}
For integers \(t\ge16\), \(F\ge1\), put
\[
 R_t(F)=\frac{(F^2+1)^t}{\binom{F^2+t}{t}}.
\]
Then
\[
 |F-R_t(F)|<\tfrac12\quad\Longleftrightarrow\quad F=t!.
\]
At the solution one also has
\[
 0<F-R_t(F)<\tfrac1{32}.
\]
\end{lemma}

\begin{proof}
Write
\[
 R_t(F)=t!\prod_{i=2}^t
       \left(1-\frac{i-1}{F^2+i}\right).
\]
Every factor increases with \(F\). Consequently
\[
 R_t(F)\ge R_t(1)=\frac{2^t}{t+1}>t^2+\tfrac12
       \qquad(t\ge16).
\]
The last inequality holds at \(16\) and persists by elementary
induction. The assumed nearest-integer bound therefore gives \(F>t^2\).
Using \(\prod_i(1-e_i)\ge1-\sum_i e_i\) for \(0\le e_i\le1\),
we obtain
\[
 R_t(F)>
 t!\left(1-\frac{t(t-1)}{2F^2}\right)>\frac34t!.
\]
Thus \(F>R_t(F)-1/2>t!/2\). Applying the same estimate a second
time, now at this much larger lower bound, gives
\[
 0<t!-R_t(F)<\frac{t!\,t(t-1)}{2F^2}
                 <\frac{2t(t-1)}{t!}<\frac1{32}.
\]
The last bound holds at \(16\), and its left-hand expression decreases
by a factor \(1/(t-1)\) when \(t\) increases by one.
It follows that
\[
 |F-t!|\le |F-R_t(F)|+|R_t(F)-t!|<\frac{17}{32}<1.
\]
Both \(F\) and \(t!\) are integers, so \(F=t!\).

Conversely, set \(F=t!\). The same product estimate directly gives
\[
 0<F-R_t(F)<\frac{t(t-1)}{2t!}<\frac1{32},
\]
which implies the nearest-integer bound.
\end{proof}

Equations \eqref{eq:near} and \eqref{eq:B-exact} provide exactly
the hypothesis of Lemma~\ref{lem:nearest-factorial}. Hence
\begin{equation}\label{eq:factorial-recovered}
 F=t!,\qquad
 B=\prod_{i=1}^t\left(1+\frac{F^2}{i}\right),\qquad
 0<BF-V<B/32.
\end{equation}
The old nonnegative factorial slacks can now be reconstructed as
\(s=BF-V\) and \(r=B-s-1\). Both are nonnegative integers.
Also \(q\ge B^{t+1}>B\), so the old output slack \(q-B-1\)
is nonnegative. Fresh \(G\)-witnesses can therefore fill each
deleted scalar certificate. Their signs are conclusions of the proof.

\section{The remaining digits and the calibrated refutation}

The beta spacing was not enlarged or otherwise modified.
From \(V=(F^2+1)^t\) one has \(F^2\mid V-1\), and therefore
\(F^2\mid y\). The standard cofinal CRT proof for
\(y=V(V-1)\) applies literally. In particular \(y>F+2\).
In the first endpoint equation, \(u_0\le-1\) would imply
\(x\le F-y<0\), contradicting \(x\ge t+1\). Thus
\(u_0\ge0\) and \(x\ge F+1\).
The second endpoint, together with \(x>1\), gives \(u_t\ge1\) and
\begin{equation}\label{eq:strong-endpoint}
 x\ge2+(t+1)y.
\end{equation}
All four range denominators in the table are now positive divisors
of \(M=q^x\).

For reference, the retained strict amplification estimates are stronger
than mere denominator divisibility. Set
\(\mathcal R=2+t+x+y+F\). Equation \eqref{eq:strong-endpoint}
and \(y>F+2\) give \(\mathcal R<2x\), while
\(x\ge F+1\ge16!+1>2^{32}\). The inherited local polynomial
ledger uses exponent \(2^{23}\), and
\[
 2^{23}\log_2(2x)<x\qquad(x\ge2^{32}).
\]
Hence its local bound is less than
\(\mathcal R^{2^{23}}<(2x)^{2^{23}}<2^x\).
This is exactly the amplification margin used in the parent
prime-power extraction, with precision \(2x\) and loss less than \(x\).

After the factorial interpretation, the encoded-index congruence is
\(B\mid C_0+F^2+1\). Since \(B>F^2+1\), the nonnegative
representative \(C_0\) cannot vanish; consequently \(H\ge2\).
For any ordinary representative \(a\ge0\) and relevant range position
\(k\le x\), use
\[
 \binom aj\le\max(1,a)^k\le H^x\qquad(0\le j\le k).
\]
The maximum is necessary when \(a=j=0\).
For the derived label, the retained degree ledger gives
\(h\le H^{2^{22}}\), and the factorial and endpoint bounds give
\(2^{22}(t+1)<x\). Its required coefficients have the same
upper bound \(H^x\).

As \(q=B^xH^x\), every true low remainder for the other ten digits
is nonnegative and satisfies
\[
 r_i^0<\frac{H^x}{q-1}q^{k_i}
       <\frac{2}{B^x}q^{k_i}<q^{k_i}/32.
\]
The shared norm and normalized equations therefore recover these
ordinary low remainders exactly. Their quotient congruences are precisely
the parent's \(N\)-divisibility range tests, with \(N=B^x\mid q\).
The inherited prime-power argument recovers the bounded local
parameters. The constructor equations and Horn-flow identities are
unchanged, so they decode a finite refutation in the numerical
\(\mathrm{C222}\) calculus \cite{common,repo}.

This is a statement about the actual nonnegative semantic
representatives. A stored representative-minus-one coordinate may
equal \(-1\); no pointwise claim that all old natural coordinates have
become nonnegative is needed or made.

The retained logical interfaces are named
\[
 \begin{gathered}
 \texttt{KEndpoint.c222\_iff\_zf\_refutation}
 \\
 \texttt{ConstructibleChoice.zf\_refutation\_iff\_zfc\_refutation}.
 \end{gathered}
\]
Their ordinary calibration identifies existence of the decoded
refutation with existence of a finite ZFC refutation in the fixed
numeration. The retained source proofs explain the underlying syntax,
rules, and calibration. This completes the soundness direction at the
arithmetic source level.

\section{Constructive completeness}

Conversely, begin with a finite calibrated refutation. The inherited
cofinal construction chooses an arbitrarily large horizon
\(t=16+d^2\), puts \(F=t!\), \(V=(F^2+1)^t\),
\(y=V(V-1)\), and supplies a sufficiently large beta-CRT code \(x\).
The two genuine beta quotients are nonnegative and less than \(x\),
since their denominators are at least two. The strict range tests
therefore admit the same canonical witnesses as the parent.

There are two parity conditions on the amplified representatives.
The factors \(1+F^2/i\) are all odd, because
\(F^2/i=F(F/i)\) is even. Hence \(B\), \(N=B^x\), and
the precision modulus \(N^2\) are odd. Changing a representative
by \(N^2\) preserves its CRT residues and reverses its parity.

Choose the dominant representative \(a\) in
\([2N^2,4N^2)\), and choose \(b\) with opposite parity.
The retained literal constructors give
\[
 h\ge2a^2,\qquad h\equiv1+a+b\pmod2.
\]
Thus \(h\) is even and remains a dominant semantic anchor.
Choose the parity of a third independent representative so that the
sum of all nine representatives is even. Each other representative
can remain in \((0,2N^2]\). These choices preserve the original
dominance and precision estimates and make \(H\) odd.
Now \(q=(BH)^x\) is odd, and
\[
 z=(q-1)/2
\]
is the required integer witness for \eqref{eq:odd}.
There is no need to prescribe the parity of \(x\).
This is weaker than the inherited odd-square construction
\cite{centered}, whose parity argument supplied an even \(x\).

For completeness of the deterministic center, take a positive
admissible \(A\), \(C=Y_h(A)\), and \(U=X_h(A)\).
Since \(h\) is even, \(C\) is even and \(U\) is odd.
Choose \(l=hC\). Pell composition gives
\[
 J_{\rm p}=Y_l(A)=C\,Y_C(U).
\]
The polynomial expansion and
\(U^2-1=(A^2-1)C^2\) give
\[
 \frac{Y_C(U)}{C}\equiv U^{C-1}\pmod C.
\]
Thus \(J_{\rm p}/C^2\) is an odd integer, supplying \(e\).
Take \(H_{\rm p}=X_l(A)\), \(f=(1-A)H_{\rm p}\), and
\(D=A+H_{\rm p}f\). The comparison pair
\((X_h(D),Y_h(D))\) has the required abscissa congruence modulo
\(H_{\rm p}\) and ordinate congruence modulo \(C\).
All quotients are integers, even if \(D\) is negative.
The retained cofinal frame and squared-cap construction supplies its
other witnesses. The full orientation pool holds because canonical
Pell coordinates have residues \(1,h\) modulo \(A-1\).

Finally \(B\mid q\mid M=q^x\). Set
\[
 v=M/B,\qquad S=v(BF-V).
\]
Lemma~\ref{lem:nearest-factorial} gives
\[
 0<S/M=F-V/B<1/32.
\]
For each genuine digit remainder choose
\(R_i=(M/L_i)r_i^0\). The preceding bounds give
\(0\le R_i<M/32\). Consequently
\[
 4\left(S^2+\sum_{i=1}^{11}R_i^2\right)
       <\frac{48}{1024}M^2<M^2.
\]
The integer
\(M^2-1-4(S^2+\sum R_i^2)\) is nonnegative and therefore
has a \(G\)-representation. This supplies the single new norm row.
All other arithmetic, constructor, and Horn rows retain their canonical
witnesses. Every finite refutation thus yields an integer zero of
the new source.

\section{From the source to the delivered quartics}

The source-level modifications above are justified by the unbounded
arithmetic proof. The subsequent reductions are exact finite algebra
and require no further number-theoretic premise \cite{residual}.
Their certificates distinguish three operations.

First, an affine integer coordinate shear or translation is a bijection
of the integer lattice, with its stated integral inverse.
Second, an elementary residual replacement
\[
 Q_i\longleftarrow Q_i+g(\mathbf X)Q_j,\qquad i\ne j,
\]
preserves the common zero set and has the explicit inverse with
\(-g\); the final coefficient reductions use integer row shears.
Third, a residual \(X_j-f(\widehat{\mathbf X})=0\) with unit
coefficient allows integral elimination of \(X_j\), with the unique
restoration \(X_j=f(\widehat{\mathbf X})\). Each actual substitution
and degree bound is checked on its literal coefficient list.

The primary endpoint uses six coordinate shears, five elementary
exposure actions, four unit-affine projections, and 409 final
unimodular residual shears after the arithmetic core.
These operations explain why dimension can fall while the final
sum of squares becomes denser and its coefficients become smaller.
An invertible row operation preserves simultaneous zeros; it need
not preserve the numerical value of the sum of squares.
The polynomial in \eqref{eq:main-definition} is the sum of squares
of the final rows, not an assertion that all intermediate quartics
are coefficientwise identical.

The supplied verifier replays the finite transformations, checks
their inverses or unit pivots, translates coordinates by the stated
convention, expands every squared residual, collects every
coefficient, and recomputes the reported metrics. The accompanying
receipts record the completed checks and rejection controls.
The source snapshot and retained proof dependencies are pinned;
the construction does not depend on a moving repository branch.

Over the integers, a sum of squares is zero exactly when every
summand is zero. Combining this fact, the certified finite
transformations, source soundness, source completeness, and the
retained logical calibration gives
\[
 \mathrm{ZFC}\ \vdash\
 \left[
   \exists\mathbf X\in\mathbb Z^{\MainVars}\;P(\mathbf X)=0
   \ \longleftrightarrow\
   \exists p\;\operatorname{Prf}_{\mathrm{ZFC}}(p,\bot)
 \right].
\]
Negation yields Theorem~\ref{thm:main}.

\section{What has and has not been established}

The completed mathematical contribution is a shared certificate for
the factorial error and digit errors, together with a signed-radix
bootstrap that removes the need for the original three initial
scalar sign graphs. The delicate points are explicit: both unit
radices are excluded before assuming a nonnegative exponent;
the first binomial digit is recovered in negative radix by absolute
remainder uniqueness; the \(F=0\) coefficients are bounded separately;
and the symmetric error condition is proved to force \(F=t!\)
by a nearest-integer theorem. The beta spacing and all fourteen
principal orientations remain intact.

The result is an ordinary mathematical proof, composed with identified
inherited arithmetic and logical proofs, and supported by complete
finite algebraic artifacts. It is formalizable in ZFC. The new
unbounded argument has not been emitted as a complete proof-assistant
kernel term or as a literal object-language ZFC derivation.
The finite checker authenticates the polynomial and its algebraic
maps; it does not prove the universal Pell, CRT, or calibration
theorems by numerical testing. Internal AI-assisted review is
documented as such and is not external peer review.

No consistency assumption is used in the equivalence, and no decision
of the equation's solvability or proof of \(\operatorname{Con}(\mathrm{ZFC})\)
is claimed. The improvement is substantial within its declared
direct quartic architecture, while its global mathematical novelty
and optimality remain separate questions. Its concrete output is
the explicit family of equations, the proof of the new arithmetic
reduction, and the reproducible coefficient and transformation data.

\begin{thebibliography}{9}\raggedright\small
\bibitem{repo}
V. Reshetnikov and ProveIt contributors, \emph{ProveIt}, inspected commit
\texttt{3ecd6a6cf53d47519228328567f45fb7ee9781d9}, 25 October 2026.
\url{https://github.com/VladimirReshetnikov/ProveIt/tree/3ecd6a6cf53d47519228328567f45fb7ee9781d9}.
\bibitem{common}
ProveIt incoming contribution, \emph{ConZFC Common Remainders Radix47},
25 October 2026. Retained 252-integer direct equation and proof closure.
\url{https://github.com/VladimirReshetnikov/ProveIt/blob/3ecd6a6cf53d47519228328567f45fb7ee9781d9/docs/incoming/ConZFC_Common_Remainders_Radix47_2026-10-25.zip}.
\bibitem{centered}
ProveIt incoming contribution, \emph{ConZFC Centered Residues},
25 October 2026. Source of the odd-square radix predecessor and the
inspected compact 75/626, 76/464, and 77/460 points.
\url{https://github.com/VladimirReshetnikov/ProveIt/blob/3ecd6a6cf53d47519228328567f45fb7ee9781d9/docs/incoming/ConZFC_Centered_Residues_2026-10-25.zip}.
\bibitem{symmetric}
ProveIt incoming contribution, \emph{ConZFC Symmetric Remainders},
25 October 2026. The deterministic-center proof is retained; its
selected-nine orientation optimization is not used here.
\url{https://github.com/VladimirReshetnikov/ProveIt/blob/3ecd6a6cf53d47519228328567f45fb7ee9781d9/docs/incoming/ConZFC_Symmetric_Remainders_2026-10-25.zip}.
\bibitem{residual}
ProveIt incoming contribution, \emph{ConZFC Residual Bases and Denominator
Engineering}, 25 October 2026. Source of the compared affine-exposure
and residual-basis methodology.
\url{https://github.com/VladimirReshetnikov/ProveIt/blob/3ecd6a6cf53d47519228328567f45fb7ee9781d9/docs/incoming/ConZFC_Residual_Bases_Denominator_Engineering_2026-10-25.zip}.
\bibitem{overview}
ProveIt, \emph{Con(ZFC) in Diophantine form}, collection overview,
and \emph{Later arrivals and current size reductions}, active review,
at the inspected commit. Counts are compared in their stated domains.
\url{https://github.com/VladimirReshetnikov/ProveIt/tree/3ecd6a6cf53d47519228328567f45fb7ee9781d9/Computability/DiophantineConsistency}.
\end{thebibliography}
\appendix
\section{The complete residual list}\label{app:residuals}
The variables \(x_0,\ldots,x_{\MainLast}\) range independently over \(\Z\).
The entire equation is \(P=\sum_{j=0}^{\MainRowLast}Q_j^2=0\).
The following list is generated from the literal JSON and independently parsed back into its exact integer coefficient data. Each monomial coefficient is literal. There are no additional quantified registers.
\begingroup\small\setlength{\jot}{1pt}
% BEGIN_Q 0
\begin{align*}
Q_{0}={}&-x_{3}^{2}+x_{4}
\end{align*}
% END_Q 0
% BEGIN_Q 1
\begin{align*}
Q_{1}={}&x_{0}-x_{0}\,x_{2}-x_{2}+x_{9}
\end{align*}
% END_Q 1
% BEGIN_Q 2
\begin{align*}
Q_{2}={}&-x_{2}\,x_{7}-x_{3}+x_{8}\,x_{9}
\end{align*}
% END_Q 2
% BEGIN_Q 3
\begin{align*}
Q_{3}={}&-1+x_{1}-x_{8}\,x_{9}
\end{align*}
% END_Q 3
% BEGIN_Q 4
\begin{align*}
Q_{4}={}&1+x_{4}-x_{5}\,x_{20}+x_{18}
\end{align*}
% END_Q 4
% BEGIN_Q 5
\begin{align*}
Q_{5}={}&-1-x_{10}^{2}-2\,x_{10}\,x_{11}-x_{11}^{2}+x_{21}
\end{align*}
% END_Q 5
% BEGIN_Q 6
\begin{align*}
Q_{6}={}&2+x_{1}-x_{3}\,x_{11}-x_{4}+x_{5}\,x_{20}-x_{8}\,x_{9}
\\ &{}+x_{10}+x_{10}^{2}-x_{10}\,x_{11}+x_{10}\,x_{12}-8\,x_{10}\,x_{21}
\\ &{}-4\,x_{10}\,x_{25}-12\,x_{11}-5\,x_{11}^{2}+7\,x_{11}\,x_{12}-16\,x_{11}\,x_{21}
\\ &{}-x_{18}-4\,x_{21}+12\,x_{21}^{2}+x_{22}+x_{25}+x_{26}+2\,x_{27}
\\ &{}-x_{40}-x_{48}+x_{78}+x_{83}
\end{align*}
% END_Q 6
% BEGIN_Q 7
\begin{align*}
Q_{7}={}&-1-x_{1}+x_{8}\,x_{9}+2\,x_{10}-x_{10}^{2}-4\,x_{10}\,x_{22}-4\,x_{22}
\\ &{}-4\,x_{22}^{2}+x_{23}
\end{align*}
% END_Q 7
% BEGIN_Q 8
\begin{align*}
Q_{8}={}&-24\,x_{10}\,x_{23}+24\,x_{10}\,x_{37}+2\,x_{23}-4\,x_{23}^{2}+x_{24}-2\,x_{37}
\\ &{}+4\,x_{37}^{2}-x_{38}
\end{align*}
% END_Q 8
% BEGIN_Q 9
\begin{align*}
Q_{9}={}&-1-x_{1}+x_{8}\,x_{9}-x_{11}-x_{11}^{2}-2\,x_{11}\,x_{12}-x_{12}
\\ &{}-x_{12}^{2}+x_{25}
\end{align*}
% END_Q 9
% BEGIN_Q 10
\begin{align*}
Q_{10}={}&-2-x_{1}+x_{3}\,x_{11}+x_{8}\,x_{9}+7\,x_{10}-x_{10}^{2}
\\ &{}+x_{10}\,x_{11}+3\,x_{10}\,x_{12}-4\,x_{10}\,x_{25}+x_{11}+2\,x_{11}^{2}
\\ &{}+3\,x_{11}\,x_{12}-x_{12}+3\,x_{12}^{2}+2\,x_{25}+x_{26}-x_{48}-x_{78}
\\ &{}-x_{83}
\end{align*}
% END_Q 10
% BEGIN_Q 11
\begin{align*}
Q_{11}={}&-1+x_{10}-x_{10}^{2}-2\,x_{10}\,x_{12}+x_{12}-x_{12}^{2}+x_{27}
\end{align*}
% END_Q 11
% BEGIN_Q 12
\begin{align*}
Q_{12}={}&1+2\,x_{0}+x_{1}-x_{3}\,x_{11}-x_{8}\,x_{9}+x_{10}+4\,x_{10}^{2}
\\ &{}-13\,x_{10}\,x_{11}+x_{10}\,x_{12}-16\,x_{10}\,x_{21}+8\,x_{10}\,x_{65}+4\,x_{11}
\\ &{}+4\,x_{11}^{2}-9\,x_{11}\,x_{12}+8\,x_{11}\,x_{25}-4\,x_{11}\,x_{27}+16\,x_{11}\,x_{65}
\\ &{}+6\,x_{12}+7\,x_{12}^{2}-2\,x_{12}\,x_{13}-24\,x_{12}\,x_{21}-x_{13}^{2}-13\,x_{21}
\\ &{}+8\,x_{21}^{2}-8\,x_{21}\,x_{27}-8\,x_{21}\,x_{65}+6\,x_{25}-4\,x_{25}^{2}
\\ &{}-8\,x_{25}\,x_{27}+2\,x_{26}+7\,x_{27}-4\,x_{27}^{2}+x_{28}+x_{30}-x_{33}
\\ &{}-8\,x_{35}-8\,x_{35}^{2}+2\,x_{36}-x_{40}+x_{43}-2\,x_{48}+4\,x_{65}
\\ &{}-4\,x_{65}^{2}+x_{66}+7\,x_{78}+x_{83}-2\,x_{214}^{2}
\end{align*}
% END_Q 12
% BEGIN_Q 13
\begin{align*}
Q_{13}={}&-4\,x_{26}^{2}-8\,x_{26}\,x_{28}-4\,x_{28}^{2}+x_{29}
\end{align*}
% END_Q 13
% BEGIN_Q 14
\begin{align*}
Q_{14}={}&-2-x_{1}+x_{3}\,x_{11}+x_{8}\,x_{9}-x_{10}-x_{10}^{2}
\\ &{}+4\,x_{10}\,x_{12}+4\,x_{10}\,x_{25}+8\,x_{10}\,x_{27}-5\,x_{11}+3\,x_{11}^{2}
\\ &{}+6\,x_{11}\,x_{12}+8\,x_{11}\,x_{25}+8\,x_{11}\,x_{27}-x_{12}+3\,x_{12}^{2}
\\ &{}-4\,x_{21}+9\,x_{25}-4\,x_{25}^{2}-8\,x_{25}\,x_{27}+x_{26}+8\,x_{27}
\\ &{}-4\,x_{27}^{2}+x_{30}-x_{48}-x_{83}
\end{align*}
% END_Q 14
% BEGIN_Q 15
\begin{align*}
Q_{15}={}&-1-x_{1}+x_{3}\,x_{11}+x_{8}\,x_{9}-7\,x_{10}+7\,x_{10}^{2}
\\ &{}+7\,x_{10}\,x_{11}-3\,x_{10}\,x_{12}+8\,x_{10}\,x_{21}+4\,x_{10}\,x_{25}
\\ &{}-8\,x_{10}\,x_{27}-8\,x_{10}\,x_{30}+2\,x_{11}-15\,x_{11}^{2}+11\,x_{11}\,x_{12}
\\ &{}-16\,x_{11}\,x_{21}+8\,x_{11}\,x_{25}-24\,x_{11}\,x_{27}-16\,x_{11}\,x_{30}-x_{12}
\\ &{}+4\,x_{12}^{2}-2\,x_{12}\,x_{13}+8\,x_{12}\,x_{27}+4\,x_{12}\,x_{30}-x_{13}^{2}
\\ &{}+2\,x_{21}-4\,x_{21}^{2}-16\,x_{21}\,x_{27}-8\,x_{21}\,x_{30}+7\,x_{25}-4\,x_{25}^{2}
\\ &{}-8\,x_{25}\,x_{27}+x_{26}+16\,x_{27}-4\,x_{27}^{2}+5\,x_{30}+x_{31}
\\ &{}+x_{43}-x_{47}-x_{48}+x_{78}-x_{83}
\end{align*}
% END_Q 15
% BEGIN_Q 16
\begin{align*}
Q_{16}={}&2+2\,x_{0}+x_{1}-x_{3}\,x_{11}-x_{4}+x_{5}\,x_{20}
\\ &{}-x_{8}\,x_{9}-2\,x_{10}-x_{10}^{2}-7\,x_{10}\,x_{11}-5\,x_{10}\,x_{12}
\\ &{}+8\,x_{10}\,x_{27}+4\,x_{10}\,x_{57}-11\,x_{11}+9\,x_{11}^{2}+x_{11}\,x_{12}
\\ &{}-8\,x_{11}\,x_{25}+8\,x_{11}\,x_{27}-16\,x_{12}-2\,x_{12}^{2}-8\,x_{12}\,x_{25}
\\ &{}+8\,x_{12}\,x_{57}-x_{18}-x_{21}-12\,x_{25}-16\,x_{25}^{2}-8\,x_{25}\,x_{27}
\\ &{}+x_{26}+7\,x_{27}-4\,x_{27}^{2}+x_{30}+x_{32}+4\,x_{35}+4\,x_{35}^{2}
\\ &{}-x_{36}+x_{48}+4\,x_{57}+4\,x_{57}^{2}-x_{58}-7\,x_{78}+x_{83}
\\ &{}-2\,x_{214}^{2}
\end{align*}
% END_Q 16
% BEGIN_Q 17
\begin{align*}
Q_{17}={}&-1+2\,x_{10}+10\,x_{10}^{2}-7\,x_{10}\,x_{11}-9\,x_{10}\,x_{12}-8\,x_{10}\,x_{25}
\\ &{}+8\,x_{10}\,x_{27}+24\,x_{10}\,x_{41}-x_{11}+6\,x_{11}^{2}+x_{11}\,x_{12}
\\ &{}-4\,x_{11}\,x_{27}-8\,x_{12}-3\,x_{12}^{2}+2\,x_{12}\,x_{13}-8\,x_{12}\,x_{27}
\\ &{}+16\,x_{12}\,x_{41}+x_{13}^{2}+4\,x_{21}-3\,x_{25}+2\,x_{26}-4\,x_{27}
\\ &{}+8\,x_{27}\,x_{41}+x_{33}+4\,x_{35}+4\,x_{35}^{2}-x_{36}-4\,x_{41}
\\ &{}+4\,x_{41}^{2}-x_{42}-x_{43}-2\,x_{48}-9\,x_{78}
\end{align*}
% END_Q 17
% BEGIN_Q 18
\begin{align*}
Q_{18}={}&1+x_{1}-x_{8}\,x_{9}-4\,x_{32}^{2}-8\,x_{32}\,x_{33}-4\,x_{33}^{2}
\\ &{}+x_{34}
\end{align*}
% END_Q 18
% BEGIN_Q 19
\begin{align*}
Q_{19}={}&1-2\,x_{10}+4\,x_{10}^{2}+4\,x_{10}\,x_{11}-4\,x_{10}\,x_{12}-8\,x_{10}\,x_{27}
\\ &{}-8\,x_{10}\,x_{65}-20\,x_{11}-7\,x_{11}^{2}+14\,x_{11}\,x_{12}+8\,x_{11}\,x_{21}
\\ &{}+16\,x_{11}\,x_{25}-8\,x_{11}\,x_{27}-16\,x_{11}\,x_{65}+x_{12}-6\,x_{12}^{2}
\\ &{}+2\,x_{12}\,x_{13}+x_{13}^{2}+7\,x_{21}-8\,x_{21}\,x_{25}+8\,x_{21}\,x_{65}
\\ &{}+2\,x_{25}+8\,x_{25}\,x_{27}-4\,x_{27}+4\,x_{27}^{2}-x_{30}+x_{35}-x_{43}
\\ &{}-4\,x_{65}+4\,x_{65}^{2}-x_{66}+6\,x_{78}
\end{align*}
% END_Q 19
% BEGIN_Q 20
\begin{align*}
Q_{20}={}&-1-x_{1}+x_{8}\,x_{9}-5\,x_{10}-x_{10}^{2}-2\,x_{10}\,x_{12}+x_{11}
\\ &{}+x_{11}^{2}+2\,x_{11}\,x_{12}-2\,x_{12}-x_{25}-x_{27}-4\,x_{35}
\\ &{}-4\,x_{35}^{2}+x_{36}
\end{align*}
% END_Q 20
% BEGIN_Q 21
\begin{align*}
Q_{21}={}&1+x_{0}+x_{1}-x_{8}\,x_{9}+11\,x_{10}+12\,x_{10}^{2}
\\ &{}+9\,x_{10}\,x_{11}-7\,x_{10}\,x_{12}-8\,x_{10}\,x_{21}+16\,x_{10}\,x_{22}
\\ &{}+12\,x_{10}\,x_{25}+8\,x_{10}\,x_{27}+8\,x_{10}\,x_{30}+6\,x_{11}-13\,x_{11}^{2}
\\ &{}+5\,x_{11}\,x_{12}-8\,x_{11}\,x_{21}+64\,x_{11}\,x_{22}+16\,x_{11}\,x_{25}
\\ &{}+24\,x_{11}\,x_{27}+16\,x_{11}\,x_{30}-15\,x_{12}+4\,x_{12}^{2}-2\,x_{12}\,x_{13}
\\ &{}+4\,x_{12}\,x_{21}+8\,x_{12}\,x_{27}-4\,x_{12}\,x_{30}-x_{13}^{2}+17\,x_{21}
\\ &{}-20\,x_{21}^{2}+32\,x_{21}\,x_{22}-8\,x_{21}\,x_{25}+16\,x_{21}\,x_{27}
\\ &{}+8\,x_{21}\,x_{30}+2\,x_{22}-16\,x_{22}^{2}+3\,x_{23}-2\,x_{24}-6\,x_{25}
\\ &{}+4\,x_{25}^{2}-7\,x_{26}+3\,x_{27}-2\,x_{29}-5\,x_{30}-3\,x_{31}-4\,x_{32}
\\ &{}-2\,x_{34}-3\,x_{35}-4\,x_{35}^{2}-x_{36}+x_{37}-4\,x_{38}-2\,x_{39}
\\ &{}+x_{40}-2\,x_{42}+x_{43}+x_{45}-2\,x_{46}+x_{47}+x_{48}
\\ &{}-x_{65}-4\,x_{67}-2\,x_{68}-3\,x_{78}-x_{210}\,x_{223}-x_{214}^{2}+x_{218}
\end{align*}
% END_Q 21
% BEGIN_Q 22
\begin{align*}
Q_{22}={}&1+x_{0}+x_{1}-x_{8}\,x_{9}-25\,x_{10}^{2}+8\,x_{10}\,x_{11}
\\ &{}+8\,x_{10}\,x_{12}+4\,x_{10}\,x_{25}-24\,x_{10}\,x_{37}+3\,x_{11}^{2}-2\,x_{11}\,x_{12}
\\ &{}-x_{12}+x_{12}^{2}-5\,x_{21}-2\,x_{25}-x_{26}-3\,x_{27}+4\,x_{35}
\\ &{}+4\,x_{35}^{2}-x_{36}+2\,x_{37}-4\,x_{37}^{2}+x_{38}+x_{48}+2\,x_{78}
\\ &{}-x_{214}^{2}
\end{align*}
% END_Q 22
% BEGIN_Q 23
\begin{align*}
Q_{23}={}&-x_{11}^{2}-4\,x_{11}\,x_{38}-4\,x_{38}^{2}+x_{39}
\end{align*}
% END_Q 23
% BEGIN_Q 24
\begin{align*}
Q_{24}={}&-2-x_{1}+x_{2}\,x_{7}+x_{3}\,x_{10}+x_{3}\,x_{11}+x_{4}
\\ &{}-x_{5}\,x_{20}-2\,x_{10}-2\,x_{10}^{2}-4\,x_{10}\,x_{11}-3\,x_{11}-x_{11}^{2}
\\ &{}+2\,x_{11}\,x_{12}+x_{12}+x_{12}^{2}+x_{18}-12\,x_{21}-16\,x_{21}^{2}
\\ &{}-x_{25}+x_{40}-x_{83}-x_{84}
\end{align*}
% END_Q 24
% BEGIN_Q 25
\begin{align*}
Q_{25}={}&1+x_{0}+x_{10}-x_{10}^{2}+5\,x_{10}\,x_{11}+5\,x_{10}\,x_{12}
\\ &{}+8\,x_{10}\,x_{27}+18\,x_{11}-13\,x_{11}^{2}-7\,x_{11}\,x_{12}-16\,x_{11}\,x_{25}
\\ &{}+8\,x_{11}\,x_{27}-24\,x_{11}\,x_{40}+7\,x_{12}+x_{12}^{2}-8\,x_{12}\,x_{25}
\\ &{}-16\,x_{12}\,x_{40}-6\,x_{21}+5\,x_{25}-12\,x_{25}^{2}-8\,x_{25}\,x_{27}
\\ &{}-8\,x_{25}\,x_{40}+2\,x_{26}+7\,x_{27}-4\,x_{27}^{2}+x_{30}-4\,x_{35}
\\ &{}-4\,x_{35}^{2}+x_{36}+4\,x_{40}-4\,x_{40}^{2}+x_{41}-x_{48}-x_{78}
\\ &{}-x_{214}^{2}
\end{align*}
% END_Q 25
% BEGIN_Q 26
\begin{align*}
Q_{26}={}&-1+x_{0}-x_{1}+x_{8}\,x_{9}-3\,x_{10}-5\,x_{10}^{2}-3\,x_{10}\,x_{11}
\\ &{}-11\,x_{10}\,x_{12}-4\,x_{10}\,x_{25}-8\,x_{10}\,x_{27}+4\,x_{10}\,x_{57}-5\,x_{11}
\\ &{}+10\,x_{11}^{2}+5\,x_{11}\,x_{12}+8\,x_{11}\,x_{25}+4\,x_{11}\,x_{27}-13\,x_{12}
\\ &{}+x_{12}^{2}-8\,x_{12}\,x_{25}-24\,x_{12}\,x_{27}+8\,x_{12}\,x_{57}-4\,x_{21}
\\ &{}+x_{25}-8\,x_{25}\,x_{27}+3\,x_{26}-9\,x_{27}-8\,x_{27}^{2}+x_{30}+x_{33}
\\ &{}+8\,x_{35}+8\,x_{35}^{2}-2\,x_{36}-4\,x_{48}+4\,x_{57}+4\,x_{57}^{2}-x_{58}
\\ &{}-5\,x_{78}-x_{214}^{2}
\end{align*}
% END_Q 26
% BEGIN_Q 27
\begin{align*}
Q_{27}={}&-x_{12}^{2}-2\,x_{12}\,x_{13}-x_{13}^{2}+x_{43}
\end{align*}
% END_Q 27
% BEGIN_Q 28
\begin{align*}
Q_{28}={}&1+x_{1}-x_{8}\,x_{9}+x_{10}^{2}-4\,x_{10}\,x_{43}+2\,x_{12}
\\ &{}-x_{12}^{2}+2\,x_{12}\,x_{13}-8\,x_{12}\,x_{43}+x_{13}^{2}-2\,x_{27}+3\,x_{43}
\\ &{}-4\,x_{43}^{2}+x_{44}
\end{align*}
% END_Q 28
% BEGIN_Q 29
\begin{align*}
Q_{29}={}&-1-x_{1}+x_{8}\,x_{9}-4\,x_{10}+5\,x_{10}^{2}+2\,x_{10}\,x_{11}
\\ &{}-6\,x_{10}\,x_{12}-8\,x_{10}\,x_{27}-4\,x_{10}\,x_{43}+5\,x_{11}+x_{11}^{2}
\\ &{}-8\,x_{11}\,x_{25}-8\,x_{11}\,x_{27}-13\,x_{12}+3\,x_{12}^{2}+16\,x_{12}\,x_{27}
\\ &{}-8\,x_{12}\,x_{43}+6\,x_{21}-11\,x_{25}+8\,x_{25}^{2}-2\,x_{26}+12\,x_{27}
\\ &{}-x_{30}+4\,x_{35}+4\,x_{35}^{2}-x_{36}+4\,x_{43}-4\,x_{43}^{2}+x_{44}
\\ &{}+x_{45}-6\,x_{78}
\end{align*}
% END_Q 29
% BEGIN_Q 30
\begin{align*}
Q_{30}={}&x_{10}\,x_{11}-x_{10}\,x_{12}-x_{11}-x_{11}\,x_{12}-3\,x_{12}
\\ &{}-2\,x_{12}\,x_{13}-x_{13}^{2}-16\,x_{27}^{2}-32\,x_{27}\,x_{43}-16\,x_{27}\,x_{44}
\\ &{}-16\,x_{27}\,x_{45}-3\,x_{43}-16\,x_{43}^{2}-16\,x_{43}\,x_{44}-16\,x_{43}\,x_{45}
\\ &{}-2\,x_{44}-4\,x_{44}^{2}-8\,x_{44}\,x_{45}-4\,x_{45}^{2}+x_{46}-x_{78}
\end{align*}
% END_Q 30
% BEGIN_Q 31
\begin{align*}
Q_{31}={}&2+x_{1}-x_{3}\,x_{11}-x_{4}+x_{5}\,x_{20}-x_{8}\,x_{9}
\\ &{}+x_{10}-12\,x_{10}^{2}+6\,x_{10}\,x_{11}+4\,x_{10}\,x_{12}-4\,x_{10}\,x_{25}
\\ &{}+24\,x_{10}\,x_{27}+16\,x_{10}\,x_{30}+4\,x_{11}-x_{11}^{2}-8\,x_{11}\,x_{12}
\\ &{}-8\,x_{11}\,x_{25}-8\,x_{11}\,x_{27}+x_{12}-4\,x_{12}^{2}+2\,x_{12}\,x_{13}
\\ &{}-8\,x_{12}\,x_{27}-4\,x_{12}\,x_{30}+x_{13}^{2}-x_{18}+x_{21}-8\,x_{25}
\\ &{}+4\,x_{25}^{2}+8\,x_{25}\,x_{27}-x_{26}-8\,x_{27}-12\,x_{27}^{2}-16\,x_{27}\,x_{30}
\\ &{}-x_{30}-4\,x_{30}^{2}-x_{43}+x_{47}+x_{48}+x_{83}
\end{align*}
% END_Q 31
% BEGIN_Q 32
\begin{align*}
Q_{32}={}&1+x_{1}-x_{8}\,x_{9}-x_{10}^{2}-2\,x_{10}\,x_{12}+x_{11}^{2}
\\ &{}+2\,x_{11}\,x_{12}-4\,x_{12}+8\,x_{12}\,x_{25}+x_{21}+6\,x_{25}-4\,x_{25}^{2}
\\ &{}+x_{48}-2\,x_{78}
\end{align*}
% END_Q 32
% BEGIN_Q 33
\begin{align*}
Q_{33}={}&1-x_{0}+x_{3}-x_{3}\,x_{10}-x_{3}\,x_{11}+x_{10}
\\ &{}+x_{10}\,x_{11}+5\,x_{10}\,x_{12}-4\,x_{10}\,x_{25}-x_{11}-x_{11}^{2}
\\ &{}-3\,x_{11}\,x_{12}+5\,x_{12}-3\,x_{12}^{2}+2\,x_{12}\,x_{13}+8\,x_{12}\,x_{25}
\\ &{}+12\,x_{12}\,x_{48}+x_{13}^{2}-x_{21}-10\,x_{25}-8\,x_{25}^{2}
\\ &{}-16\,x_{25}\,x_{48}+x_{26}-2\,x_{27}-4\,x_{35}-4\,x_{35}^{2}+x_{36}
\\ &{}-x_{43}-3\,x_{48}-4\,x_{48}^{2}+x_{49}+x_{78}+x_{83}+x_{84}
\\ &{}+x_{214}^{2}
\end{align*}
% END_Q 33
% BEGIN_Q 34
\begin{align*}
Q_{34}={}&1+x_{1}-x_{3}\,x_{11}-x_{8}\,x_{9}+2\,x_{11}-2\,x_{11}\,x_{12}
\\ &{}-4\,x_{11}\,x_{49}+x_{12}-3\,x_{12}^{2}-8\,x_{12}\,x_{49}-x_{25}-4\,x_{49}^{2}
\\ &{}+x_{50}+x_{83}
\end{align*}
% END_Q 34
% BEGIN_Q 35
\begin{align*}
Q_{35}={}&-1-4\,x_{50}-4\,x_{50}^{2}+x_{51}
\end{align*}
% END_Q 35
% BEGIN_Q 36
\begin{align*}
Q_{36}={}&1-x_{12}^{2}-8\,x_{12}\,x_{50}-4\,x_{12}\,x_{51}+8\,x_{50}-8\,x_{50}^{2}
\\ &{}-16\,x_{50}\,x_{51}-2\,x_{51}-4\,x_{51}^{2}+x_{52}
\end{align*}
% END_Q 36
% BEGIN_Q 37
\begin{align*}
Q_{37}={}&2+x_{1}-x_{3}\,x_{11}-x_{8}\,x_{9}+2\,x_{11}-2\,x_{11}\,x_{12}
\\ &{}-4\,x_{11}\,x_{52}+x_{12}-3\,x_{12}^{2}-8\,x_{12}\,x_{52}-x_{25}-4\,x_{52}^{2}
\\ &{}+x_{53}+x_{83}
\end{align*}
% END_Q 37
% BEGIN_Q 38
\begin{align*}
Q_{38}={}&-4\,x_{47}^{2}-8\,x_{47}\,x_{53}-4\,x_{53}^{2}+x_{54}
\end{align*}
% END_Q 38
% BEGIN_Q 39
\begin{align*}
Q_{39}={}&-1-x_{1}+x_{8}\,x_{9}-x_{10}+x_{10}^{2}-x_{10}\,x_{11}
\\ &{}-x_{10}\,x_{12}+4\,x_{10}\,x_{25}-8\,x_{10}\,x_{47}-4\,x_{10}\,x_{54}-x_{11}^{2}
\\ &{}-x_{11}\,x_{12}+x_{12}-x_{12}^{2}-3\,x_{25}-x_{26}-x_{27}
\\ &{}+24\,x_{47}-12\,x_{47}^{2}+8\,x_{47}\,x_{53}-16\,x_{47}\,x_{54}+x_{48}
\\ &{}+4\,x_{53}^{2}+11\,x_{54}-4\,x_{54}^{2}+x_{55}+x_{78}
\end{align*}
% END_Q 39
% BEGIN_Q 40
\begin{align*}
Q_{40}={}&-1-4\,x_{55}-4\,x_{55}^{2}+x_{56}
\end{align*}
% END_Q 40
% BEGIN_Q 41
\begin{align*}
Q_{41}={}&1+x_{1}-x_{8}\,x_{9}-x_{12}^{2}-8\,x_{12}\,x_{55}-4\,x_{12}\,x_{56}
\\ &{}+8\,x_{55}-8\,x_{55}^{2}-16\,x_{55}\,x_{56}-2\,x_{56}-4\,x_{56}^{2}+x_{57}
\end{align*}
% END_Q 41
% BEGIN_Q 42
\begin{align*}
Q_{42}={}&1+x_{1}-x_{8}\,x_{9}-2\,x_{10}-x_{10}^{2}-x_{10}\,x_{11}
\\ &{}-3\,x_{10}\,x_{12}-4\,x_{10}\,x_{57}+x_{11}\,x_{12}-3\,x_{12}-x_{12}^{2}
\\ &{}-8\,x_{12}\,x_{57}+x_{21}-x_{25}-x_{27}-4\,x_{57}-4\,x_{57}^{2}+x_{58}
\\ &{}-x_{78}
\end{align*}
% END_Q 42
% BEGIN_Q 43
\begin{align*}
Q_{43}={}&2-x_{0}+x_{1}-x_{2}\,x_{7}-x_{3}\,x_{10}-x_{3}\,x_{11}
\\ &{}-x_{4}+x_{5}\,x_{20}-x_{18}-12\,x_{58}-4\,x_{58}^{2}+x_{59}
\\ &{}+x_{83}+x_{84}+x_{214}^{2}
\end{align*}
% END_Q 43
% BEGIN_Q 44
\begin{align*}
Q_{44}={}&-x_{10}\,x_{11}+x_{10}\,x_{12}-x_{11}^{2}-x_{11}\,x_{12}+2\,x_{12}
\\ &{}-2\,x_{12}^{2}+16\,x_{12}\,x_{58}+8\,x_{12}\,x_{59}-3\,x_{25}-16\,x_{25}\,x_{58}
\\ &{}-8\,x_{25}\,x_{59}-x_{48}+8\,x_{58}-16\,x_{58}^{2}-16\,x_{58}\,x_{59}+4\,x_{59}
\\ &{}-4\,x_{59}^{2}+x_{60}+x_{78}
\end{align*}
% END_Q 44
% BEGIN_Q 45
\begin{align*}
Q_{45}={}&-1-x_{1}+x_{8}\,x_{9}-x_{10}+2\,x_{10}^{2}+2\,x_{10}\,x_{11}
\\ &{}+2\,x_{10}\,x_{12}-2\,x_{11}-3\,x_{11}^{2}+2\,x_{11}\,x_{12}-8\,x_{11}\,x_{25}
\\ &{}-4\,x_{11}\,x_{60}+4\,x_{12}-4\,x_{12}^{2}+2\,x_{12}\,x_{13}-12\,x_{12}\,x_{48}
\\ &{}+16\,x_{12}\,x_{60}+x_{13}^{2}-x_{21}-5\,x_{25}+4\,x_{25}^{2}+16\,x_{25}\,x_{48}
\\ &{}-16\,x_{25}\,x_{60}-x_{27}-x_{43}-x_{48}+4\,x_{48}^{2}-x_{49}
\\ &{}-4\,x_{60}^{2}+x_{61}
\end{align*}
% END_Q 45
% BEGIN_Q 46
\begin{align*}
Q_{46}={}&2-x_{0}+x_{1}-x_{2}\,x_{7}-x_{3}\,x_{10}-x_{3}\,x_{11}
\\ &{}-x_{4}+x_{5}\,x_{20}-x_{18}-12\,x_{53}-4\,x_{53}^{2}+x_{62}
\\ &{}+x_{83}+x_{84}+x_{214}^{2}
\end{align*}
% END_Q 46
% BEGIN_Q 47
\begin{align*}
Q_{47}={}&1+4\,x_{47}^{2}+8\,x_{47}\,x_{53}-12\,x_{53}^{2}-16\,x_{53}\,x_{61}
\\ &{}-16\,x_{53}\,x_{62}-x_{54}-4\,x_{61}^{2}-8\,x_{61}\,x_{62}-4\,x_{62}^{2}+x_{63}
\end{align*}
% END_Q 47
% BEGIN_Q 48
\begin{align*}
Q_{48}={}&x_{12}-x_{12}^{2}-8\,x_{12}\,x_{61}-4\,x_{12}\,x_{63}-16\,x_{61}^{2}
\\ &{}-16\,x_{61}\,x_{63}-4\,x_{63}^{2}+x_{64}
\end{align*}
% END_Q 48
% BEGIN_Q 49
\begin{align*}
Q_{49}={}&2+x_{1}-x_{3}\,x_{11}-x_{8}\,x_{9}+4\,x_{10}+2\,x_{10}^{2}
\\ &{}+8\,x_{10}\,x_{21}+3\,x_{11}-11\,x_{11}^{2}+2\,x_{11}\,x_{12}+16\,x_{11}\,x_{21}
\\ &{}+2\,x_{12}+2\,x_{12}^{2}-4\,x_{12}\,x_{21}-10\,x_{21}-4\,x_{21}^{2}+x_{25}
\\ &{}+x_{65}-4\,x_{78}+x_{83}
\end{align*}
% END_Q 49
% BEGIN_Q 50
\begin{align*}
Q_{50}={}&2\,x_{10}\,x_{12}-4\,x_{10}\,x_{25}+8\,x_{10}\,x_{65}-9\,x_{11}+5\,x_{11}^{2}
\\ &{}+2\,x_{11}\,x_{12}+16\,x_{11}\,x_{65}-3\,x_{12}-x_{12}^{2}+4\,x_{12}\,x_{21}
\\ &{}+6\,x_{21}-8\,x_{21}\,x_{65}+x_{25}+x_{26}+x_{27}-x_{48}+3\,x_{65}
\\ &{}-4\,x_{65}^{2}+x_{66}+4\,x_{78}
\end{align*}
% END_Q 50
% BEGIN_Q 51
\begin{align*}
Q_{51}={}&1+x_{1}-x_{8}\,x_{9}+2\,x_{10}+x_{10}^{2}-3\,x_{10}\,x_{11}
\\ &{}-3\,x_{10}\,x_{12}+8\,x_{10}\,x_{21}-8\,x_{10}\,x_{25}+16\,x_{10}\,x_{64}
\\ &{}+16\,x_{10}\,x_{66}-3\,x_{11}+x_{11}^{2}-3\,x_{11}\,x_{12}+16\,x_{11}\,x_{21}
\\ &{}+32\,x_{11}\,x_{64}+32\,x_{11}\,x_{66}-4\,x_{12}+3\,x_{12}^{2}-2\,x_{12}\,x_{13}
\\ &{}+12\,x_{12}\,x_{21}-x_{13}^{2}-x_{21}-4\,x_{21}^{2}-16\,x_{21}\,x_{64}
\\ &{}-16\,x_{21}\,x_{66}+2\,x_{26}-4\,x_{35}-4\,x_{35}^{2}+x_{36}+x_{43}
\\ &{}-2\,x_{48}-4\,x_{64}^{2}-8\,x_{64}\,x_{66}-3\,x_{65}-2\,x_{66}-4\,x_{66}^{2}
\\ &{}+x_{67}+5\,x_{78}
\end{align*}
% END_Q 51
% BEGIN_Q 52
\begin{align*}
Q_{52}={}&-1-x_{1}+x_{8}\,x_{9}-4\,x_{67}-4\,x_{67}^{2}+x_{68}
\end{align*}
% END_Q 52
% BEGIN_Q 53
\begin{align*}
Q_{53}={}&-x_{3}\,x_{15}+x_{69}
\end{align*}
% END_Q 53
% BEGIN_Q 54
\begin{align*}
Q_{54}={}&x_{14}-x_{15}-2\,x_{24}\,x_{69}-x_{70}
\end{align*}
% END_Q 54
% BEGIN_Q 55
\begin{align*}
Q_{55}={}&x_{14}-x_{15}-4\,x_{26}\,x_{69}-2\,x_{29}\,x_{69}-x_{71}
\end{align*}
% END_Q 55
% BEGIN_Q 56
\begin{align*}
Q_{56}={}&-x_{3}\,x_{15}-2\,x_{31}\,x_{69}+2\,x_{46}\,x_{69}-x_{69}-x_{72}
\\ &{}+x_{77}
\end{align*}
% END_Q 56
% BEGIN_Q 57
\begin{align*}
Q_{57}={}&x_{14}-x_{15}-4\,x_{32}\,x_{69}-2\,x_{34}\,x_{69}-x_{73}
\end{align*}
% END_Q 57
% BEGIN_Q 58
\begin{align*}
Q_{58}={}&x_{14}-x_{15}+2\,x_{24}\,x_{69}-2\,x_{31}\,x_{69}-2\,x_{36}\,x_{69}
\\ &{}+x_{70}-x_{72}-x_{74}
\end{align*}
% END_Q 58
% BEGIN_Q 59
\begin{align*}
Q_{59}={}&x_{14}-x_{15}-4\,x_{38}\,x_{69}-2\,x_{39}\,x_{69}-x_{75}
\end{align*}
% END_Q 59
% BEGIN_Q 60
\begin{align*}
Q_{60}={}&x_{14}-x_{15}-2\,x_{42}\,x_{69}-x_{76}
\end{align*}
% END_Q 60
% BEGIN_Q 61
\begin{align*}
Q_{61}={}&-2\,x_{3}\,x_{15}+x_{14}-x_{15}-2\,x_{46}\,x_{69}-2\,x_{69}-x_{77}
\end{align*}
% END_Q 61
% BEGIN_Q 62
\begin{align*}
Q_{62}={}&-x_{10}\,x_{11}+x_{10}\,x_{12}+x_{11}+x_{11}\,x_{12}-x_{12}
\\ &{}-x_{12}^{2}+x_{78}
\end{align*}
% END_Q 62
% BEGIN_Q 63
\begin{align*}
Q_{63}={}&1+2\,x_{13}+x_{13}^{2}-x_{78}^{2}+x_{79}
\end{align*}
% END_Q 63
% BEGIN_Q 64
\begin{align*}
Q_{64}={}&-x_{79}^{2}+x_{80}
\end{align*}
% END_Q 64
% BEGIN_Q 65
\begin{align*}
Q_{65}={}&x_{3}\,x_{15}-4\,x_{67}\,x_{69}-2\,x_{68}\,x_{69}+x_{69}+x_{81}
\end{align*}
% END_Q 65
% BEGIN_Q 66
\begin{align*}
Q_{66}={}&-x_{14}^{2}+2\,x_{14}\,x_{15}+2\,x_{14}\,x_{81}-x_{15}^{2}-2\,x_{15}\,x_{81}
\\ &{}-x_{81}^{2}+x_{82}
\end{align*}
% END_Q 66
% BEGIN_Q 67
\begin{align*}
Q_{67}={}&-1-x_{3}\,x_{11}+x_{83}
\end{align*}
% END_Q 67
% BEGIN_Q 68
\begin{align*}
Q_{68}={}&-x_{15}\,x_{83}+x_{85}\,x_{86}+x_{87}
\end{align*}
% END_Q 68
% BEGIN_Q 69
\begin{align*}
Q_{69}={}&-1+x_{3}-x_{3}\,x_{10}+x_{84}
\end{align*}
% END_Q 69
% BEGIN_Q 70
\begin{align*}
Q_{70}={}&-2-x_{2}\,x_{7}+x_{3}+2\,x_{3}\,x_{10}+3\,x_{3}\,x_{11}
\\ &{}-2\,x_{3}\,x_{21}+x_{8}\,x_{9}+x_{83}+x_{85}
\end{align*}
% END_Q 70
% BEGIN_Q 71
\begin{align*}
Q_{71}={}&-x_{14}\,x_{84}+x_{86}
\end{align*}
% END_Q 71
% BEGIN_Q 72
\begin{align*}
Q_{72}={}&3\,x_{14}+x_{14}\,x_{83}+9\,x_{15}-5\,x_{15}\,x_{83}-4\,x_{15}\,x_{84}
\\ &{}-x_{15}\,x_{85}-2\,x_{24}\,x_{69}-2\,x_{36}\,x_{69}-2\,x_{42}\,x_{69}-x_{70}
\\ &{}-x_{74}-x_{76}-3\,x_{85}\,x_{86}-3\,x_{87}+x_{88}
\end{align*}
% END_Q 72
% BEGIN_Q 73
\begin{align*}
Q_{73}={}&-x_{14}\,x_{70}+x_{15}\,x_{70}+x_{89}
\end{align*}
% END_Q 73
% BEGIN_Q 74
\begin{align*}
Q_{74}={}&-x_{71}\,x_{89}+x_{90}
\end{align*}
% END_Q 74
% BEGIN_Q 75
\begin{align*}
Q_{75}={}&-x_{72}\,x_{90}+x_{91}
\end{align*}
% END_Q 75
% BEGIN_Q 76
\begin{align*}
Q_{76}={}&-x_{73}\,x_{91}+x_{92}
\end{align*}
% END_Q 76
% BEGIN_Q 77
\begin{align*}
Q_{77}={}&-x_{74}\,x_{92}+x_{93}
\end{align*}
% END_Q 77
% BEGIN_Q 78
\begin{align*}
Q_{78}={}&-x_{75}\,x_{93}+x_{94}
\end{align*}
% END_Q 78
% BEGIN_Q 79
\begin{align*}
Q_{79}={}&-x_{76}\,x_{94}+x_{95}
\end{align*}
% END_Q 79
% BEGIN_Q 80
\begin{align*}
Q_{80}={}&-x_{77}\,x_{95}+x_{96}
\end{align*}
% END_Q 80
% BEGIN_Q 81
\begin{align*}
Q_{81}={}&-x_{87}\,x_{96}+x_{97}
\end{align*}
% END_Q 81
% BEGIN_Q 82
\begin{align*}
Q_{82}={}&-x_{88}\,x_{97}+x_{98}
\end{align*}
% END_Q 82
% BEGIN_Q 83
\begin{align*}
Q_{83}={}&-x_{98}^{2}+x_{99}
\end{align*}
% END_Q 83
% BEGIN_Q 84
\begin{align*}
Q_{84}={}&-x_{2}-x_{2}\,x_{18}+x_{18}+x_{100}
\end{align*}
% END_Q 84
% BEGIN_Q 85
\begin{align*}
Q_{85}={}&-x_{16}\,x_{100}+x_{101}
\end{align*}
% END_Q 85
% BEGIN_Q 86
\begin{align*}
Q_{86}={}&-2\,x_{1}^{2}+2\,x_{1}\,x_{14}+2\,x_{1}\,x_{15}+2\,x_{1}\,x_{101}+2\,x_{1}\,x_{102}
\\ &{}-x_{14}^{2}-2\,x_{14}\,x_{101}-x_{15}^{2}-2\,x_{15}\,x_{102}+x_{19}\,x_{103}
\\ &{}-x_{80}\,x_{99}-x_{82}\,x_{99}-x_{101}^{2}-x_{102}^{2}
\end{align*}
% END_Q 86
% BEGIN_Q 87
\begin{align*}
Q_{87}={}&-x_{2}\,x_{17}+x_{17}-x_{17}\,x_{100}+x_{102}
\end{align*}
% END_Q 87
% BEGIN_Q 88
\begin{align*}
Q_{88}={}&-x_{106}-2\,x_{106}\,x_{237}+x_{108}
\end{align*}
% END_Q 88
% BEGIN_Q 89
\begin{align*}
Q_{89}={}&-x_{5}\,x_{130}+x_{104}\,x_{228}-x_{107}-x_{108}\,x_{130}
\end{align*}
% END_Q 89
% BEGIN_Q 90
\begin{align*}
Q_{90}={}&x_{109}\,x_{229}-x_{111}\,x_{133}-x_{112}
\end{align*}
% END_Q 90
% BEGIN_Q 91
\begin{align*}
Q_{91}={}&x_{113}\,x_{229}-x_{114}\,x_{133}-x_{115}
\end{align*}
% END_Q 91
% BEGIN_Q 92
\begin{align*}
Q_{92}={}&x_{116}\,x_{229}-x_{117}\,x_{133}-x_{118}
\end{align*}
% END_Q 92
% BEGIN_Q 93
\begin{align*}
Q_{93}={}&x_{119}\,x_{229}-x_{120}\,x_{133}-x_{121}
\end{align*}
% END_Q 93
% BEGIN_Q 94
\begin{align*}
Q_{94}={}&x_{122}\,x_{230}-x_{124}\,x_{133}-x_{125}
\end{align*}
% END_Q 94
% BEGIN_Q 95
\begin{align*}
Q_{95}={}&x_{126}\,x_{230}-x_{127}\,x_{133}-x_{128}
\end{align*}
% END_Q 95
% BEGIN_Q 96
\begin{align*}
Q_{96}={}&-x_{19}\,x_{130}+x_{133}
\end{align*}
% END_Q 96
% BEGIN_Q 97
\begin{align*}
Q_{97}={}&x_{129}-x_{131}\,x_{133}-x_{132}
\end{align*}
% END_Q 97
% BEGIN_Q 98
\begin{align*}
Q_{98}={}&-x_{133}\,x_{135}+x_{134}-x_{136}
\end{align*}
% END_Q 98
% BEGIN_Q 99
\begin{align*}
Q_{99}={}&-x_{133}\,x_{138}+x_{137}\,x_{228}-x_{139}
\end{align*}
% END_Q 99
% BEGIN_Q 100
\begin{align*}
Q_{100}={}&-x_{133}\,x_{141}+x_{140}\,x_{229}-x_{142}
\end{align*}
% END_Q 100
% BEGIN_Q 101
\begin{align*}
Q_{101}={}&x_{143}-2\,x_{143}\,x_{145}+x_{146}
\end{align*}
% END_Q 101
% BEGIN_Q 102
\begin{align*}
Q_{102}={}&-x_{143}\,x_{146}-x_{146}+x_{147}
\end{align*}
% END_Q 102
% BEGIN_Q 103
\begin{align*}
Q_{103}={}&x_{146}^{2}-x_{147}^{2}+x_{167}
\end{align*}
% END_Q 103
% BEGIN_Q 104
\begin{align*}
Q_{104}={}&-x_{144}\,x_{149}-x_{149}+x_{150}
\end{align*}
% END_Q 104
% BEGIN_Q 105
\begin{align*}
Q_{105}={}&2\,x_{148}+x_{148}^{2}+x_{149}^{2}-x_{150}^{2}
\end{align*}
% END_Q 105
% BEGIN_Q 106
\begin{align*}
Q_{106}={}&x_{4}^{2}-2\,x_{4}\,x_{144}+x_{144}^{2}+x_{151}-x_{167}
\end{align*}
% END_Q 106
% BEGIN_Q 107
\begin{align*}
Q_{107}={}&-1-x_{4}\,x_{149}+x_{6}+x_{144}\,x_{149}-x_{148}+x_{151}\,x_{152}
\end{align*}
% END_Q 107
% BEGIN_Q 108
\begin{align*}
Q_{108}={}&1-2\,x_{5}+x_{5}^{2}-2\,x_{5}\,x_{144}+x_{153}
\end{align*}
% END_Q 108
% BEGIN_Q 109
\begin{align*}
Q_{109}={}&-x_{5}\,x_{198}+x_{19}+x_{153}\,x_{154}-x_{197}+x_{199}
\end{align*}
% END_Q 109
% BEGIN_Q 110
\begin{align*}
Q_{110}={}&-1+2\,x_{144}\,x_{213}-x_{155}+2\,x_{213}-x_{213}^{2}
\end{align*}
% END_Q 110
% BEGIN_Q 111
\begin{align*}
Q_{111}={}&1+x_{155}\,x_{156}-x_{197}-x_{198}\,x_{213}+x_{199}+2\,x_{237}
\end{align*}
% END_Q 111
% BEGIN_Q 112
\begin{align*}
Q_{112}={}&-x_{144}\,x_{158}-x_{158}+x_{159}
\end{align*}
% END_Q 112
% BEGIN_Q 113
\begin{align*}
Q_{113}={}&2\,x_{157}+x_{157}^{2}+x_{158}^{2}-x_{159}^{2}
\end{align*}
% END_Q 113
% BEGIN_Q 114
\begin{align*}
Q_{114}={}&-x_{144}+x_{144}^{2}-4\,x_{144}\,x_{237}-x_{155}\,x_{156}+x_{160}-x_{167}
\\ &{}+x_{197}+x_{198}\,x_{213}-x_{199}+x_{215}\,x_{221}-x_{219}+2\,x_{237}
\\ &{}+4\,x_{237}^{2}
\end{align*}
% END_Q 114
% BEGIN_Q 115
\begin{align*}
Q_{115}={}&-1+x_{104}+x_{144}\,x_{158}-x_{157}-x_{158}-2\,x_{158}\,x_{237}
\\ &{}+x_{160}\,x_{161}
\end{align*}
% END_Q 115
% BEGIN_Q 116
\begin{align*}
Q_{116}={}&-1+x_{105}+x_{144}\,x_{149}-2\,x_{144}\,x_{162}-x_{148}-2\,x_{149}\,x_{237}
\\ &{}+x_{160}\,x_{162}+x_{162}+4\,x_{162}\,x_{237}
\end{align*}
% END_Q 116
% BEGIN_Q 117
\begin{align*}
Q_{117}={}&1+x_{144}-x_{144}\,x_{164}-x_{164}+x_{165}
\end{align*}
% END_Q 117
% BEGIN_Q 118
\begin{align*}
Q_{118}={}&-x_{144}+x_{144}\,x_{164}+2\,x_{163}+x_{163}^{2}-x_{164}+x_{164}^{2}
\\ &{}-x_{165}-x_{165}^{2}
\end{align*}
% END_Q 118
% BEGIN_Q 119
\begin{align*}
Q_{119}={}&x_{109}-x_{144}+x_{144}\,x_{164}+x_{160}\,x_{166}-x_{163}-x_{164}
\\ &{}-2\,x_{164}\,x_{237}+2\,x_{237}
\end{align*}
% END_Q 119
% BEGIN_Q 120
\begin{align*}
Q_{120}={}&-x_{144}-x_{144}^{2}+x_{167}+x_{215}\,x_{221}-x_{219}
\end{align*}
% END_Q 120
% BEGIN_Q 121
\begin{align*}
Q_{121}={}&x_{110}+x_{144}\,x_{150}-2\,x_{144}\,x_{168}-x_{148}-2\,x_{148}\,x_{237}
\\ &{}-x_{149}\,x_{167}-2\,x_{150}\,x_{237}+x_{155}\,x_{156}+x_{160}\,x_{168}+x_{168}
\\ &{}+4\,x_{168}\,x_{237}-x_{197}-x_{198}\,x_{213}+x_{199}
\end{align*}
% END_Q 121
% BEGIN_Q 122
\begin{align*}
Q_{122}={}&-x_{144}\,x_{170}-x_{170}+x_{171}
\end{align*}
% END_Q 122
% BEGIN_Q 123
\begin{align*}
Q_{123}={}&2\,x_{169}+x_{169}^{2}+x_{170}^{2}-x_{171}^{2}
\end{align*}
% END_Q 123
% BEGIN_Q 124
\begin{align*}
Q_{124}={}&-1+x_{113}+x_{144}\,x_{170}+x_{160}\,x_{172}-x_{169}-x_{170}
\\ &{}-2\,x_{170}\,x_{237}
\end{align*}
% END_Q 124
% BEGIN_Q 125
\begin{align*}
Q_{125}={}&-x_{144}\,x_{174}-x_{174}+x_{175}
\end{align*}
% END_Q 125
% BEGIN_Q 126
\begin{align*}
Q_{126}={}&2\,x_{173}+x_{173}^{2}+x_{174}^{2}-x_{175}^{2}
\end{align*}
% END_Q 126
% BEGIN_Q 127
\begin{align*}
Q_{127}={}&-1+x_{116}+x_{144}\,x_{174}+x_{160}\,x_{176}-x_{173}-x_{174}
\\ &{}-2\,x_{174}\,x_{237}
\end{align*}
% END_Q 127
% BEGIN_Q 128
\begin{align*}
Q_{128}={}&-x_{144}\,x_{178}-x_{178}+x_{179}
\end{align*}
% END_Q 128
% BEGIN_Q 129
\begin{align*}
Q_{129}={}&2\,x_{177}+x_{177}^{2}+x_{178}^{2}-x_{179}^{2}
\end{align*}
% END_Q 129
% BEGIN_Q 130
\begin{align*}
Q_{130}={}&-1+x_{119}+x_{144}\,x_{178}+x_{160}\,x_{180}-x_{177}-x_{178}
\\ &{}-2\,x_{178}\,x_{237}
\end{align*}
% END_Q 130
% BEGIN_Q 131
\begin{align*}
Q_{131}={}&-x_{144}\,x_{182}-x_{182}+x_{183}
\end{align*}
% END_Q 131
% BEGIN_Q 132
\begin{align*}
Q_{132}={}&2\,x_{181}+x_{181}^{2}+x_{182}^{2}-x_{183}^{2}
\end{align*}
% END_Q 132
% BEGIN_Q 133
\begin{align*}
Q_{133}={}&-1+x_{122}+x_{144}\,x_{182}+x_{160}\,x_{184}-x_{181}-x_{182}
\\ &{}-2\,x_{182}\,x_{237}
\end{align*}
% END_Q 133
% BEGIN_Q 134
\begin{align*}
Q_{134}={}&-x_{144}\,x_{186}-x_{186}+x_{187}
\end{align*}
% END_Q 134
% BEGIN_Q 135
\begin{align*}
Q_{135}={}&2\,x_{185}+x_{185}^{2}+x_{186}^{2}-x_{187}^{2}
\end{align*}
% END_Q 135
% BEGIN_Q 136
\begin{align*}
Q_{136}={}&-1+x_{123}+x_{144}\,x_{186}-2\,x_{144}\,x_{188}+x_{160}\,x_{188}-x_{185}
\\ &{}-2\,x_{186}\,x_{237}+x_{188}+4\,x_{188}\,x_{237}
\end{align*}
% END_Q 136
% BEGIN_Q 137
\begin{align*}
Q_{137}={}&-x_{144}\,x_{190}-x_{190}+x_{191}
\end{align*}
% END_Q 137
% BEGIN_Q 138
\begin{align*}
Q_{138}={}&2\,x_{189}+x_{189}^{2}+x_{190}^{2}-x_{191}^{2}
\end{align*}
% END_Q 138
% BEGIN_Q 139
\begin{align*}
Q_{139}={}&-1+x_{126}+x_{144}\,x_{190}+x_{160}\,x_{192}-x_{189}-x_{190}
\\ &{}-2\,x_{190}\,x_{237}
\end{align*}
% END_Q 139
% BEGIN_Q 140
\begin{align*}
Q_{140}={}&-x_{144}\,x_{194}-x_{194}+x_{195}
\end{align*}
% END_Q 140
% BEGIN_Q 141
\begin{align*}
Q_{141}={}&2\,x_{193}+x_{193}^{2}+x_{194}^{2}-x_{195}^{2}
\end{align*}
% END_Q 141
% BEGIN_Q 142
\begin{align*}
Q_{142}={}&-1+x_{129}+x_{144}\,x_{194}+x_{160}\,x_{196}-x_{193}-x_{194}
\\ &{}-2\,x_{194}\,x_{237}
\end{align*}
% END_Q 142
% BEGIN_Q 143
\begin{align*}
Q_{143}={}&1-x_{144}\,x_{198}-x_{198}+x_{199}
\end{align*}
% END_Q 143
% BEGIN_Q 144
\begin{align*}
Q_{144}={}&-1-x_{5}\,x_{198}+x_{19}+x_{153}\,x_{154}+x_{197}+x_{197}^{2}
\\ &{}+x_{198}^{2}-x_{199}-x_{199}^{2}
\end{align*}
% END_Q 144
% BEGIN_Q 145
\begin{align*}
Q_{145}={}&x_{130}-2\,x_{144}\,x_{200}+x_{160}\,x_{200}-x_{197}-x_{198}
\\ &{}-2\,x_{198}\,x_{237}+x_{199}+x_{200}+4\,x_{200}\,x_{237}
\end{align*}
% END_Q 145
% BEGIN_Q 146
\begin{align*}
Q_{146}={}&-x_{144}\,x_{202}-x_{202}+x_{203}
\end{align*}
% END_Q 146
% BEGIN_Q 147
\begin{align*}
Q_{147}={}&2\,x_{201}+x_{201}^{2}+x_{202}^{2}-x_{203}^{2}
\end{align*}
% END_Q 147
% BEGIN_Q 148
\begin{align*}
Q_{148}={}&-1+x_{134}+x_{144}\,x_{202}+x_{160}\,x_{204}-x_{201}-x_{202}
\\ &{}-2\,x_{202}\,x_{237}
\end{align*}
% END_Q 148
% BEGIN_Q 149
\begin{align*}
Q_{149}={}&-x_{144}\,x_{206}-x_{206}+x_{207}
\end{align*}
% END_Q 149
% BEGIN_Q 150
\begin{align*}
Q_{150}={}&2\,x_{205}+x_{205}^{2}+x_{206}^{2}-x_{207}^{2}
\end{align*}
% END_Q 150
% BEGIN_Q 151
\begin{align*}
Q_{151}={}&-1+x_{137}+x_{144}\,x_{206}+x_{160}\,x_{208}-x_{205}-x_{206}
\\ &{}-2\,x_{206}\,x_{237}
\end{align*}
% END_Q 151
% BEGIN_Q 152
\begin{align*}
Q_{152}={}&-x_{144}\,x_{210}-x_{210}+x_{211}
\end{align*}
% END_Q 152
% BEGIN_Q 153
\begin{align*}
Q_{153}={}&2\,x_{209}+x_{209}^{2}+x_{210}^{2}-x_{211}^{2}
\end{align*}
% END_Q 153
% BEGIN_Q 154
\begin{align*}
Q_{154}={}&-1+x_{140}+x_{144}\,x_{210}+x_{160}\,x_{212}-x_{209}-x_{210}
\\ &{}-2\,x_{210}\,x_{237}
\end{align*}
% END_Q 154
% BEGIN_Q 155
\begin{align*}
Q_{155}={}&-1-2\,x_{0}-3\,x_{0}^{2}-2\,x_{0}\,x_{4}-x_{1}-x_{1}^{2}-x_{2}
\\ &{}-x_{2}^{2}-x_{4}-2\,x_{4}^{2}-x_{5}^{2}-x_{5}\,x_{20}+x_{5}\,x_{198}
\\ &{}-x_{6}+x_{8}\,x_{9}+x_{10}+2\,x_{10}\,x_{12}-x_{11}^{2}-x_{12}
\\ &{}-x_{13}^{2}-x_{14}^{2}-x_{15}^{2}-x_{16}^{2}-x_{17}^{2}+x_{18}
\\ &{}-x_{18}^{2}-x_{19}-x_{19}^{2}-x_{27}-x_{104}^{2}-x_{105}^{2}
\\ &{}-x_{109}^{2}-x_{110}^{2}-x_{113}^{2}-x_{116}^{2}-x_{119}^{2}-x_{122}^{2}
\\ &{}-x_{123}^{2}-x_{126}^{2}-x_{129}^{2}-x_{130}^{2}-x_{134}^{2}-x_{137}^{2}
\\ &{}-x_{140}^{2}+x_{143}-x_{144}+x_{144}^{2}-4\,x_{144}\,x_{237}+x_{149}^{2}
\\ &{}-x_{153}\,x_{154}+2\,x_{155}\,x_{156}+x_{158}^{2}+x_{160}-2\,x_{164}+x_{164}^{2}
\\ &{}-x_{167}+x_{170}^{2}+x_{174}^{2}+x_{178}^{2}+x_{182}^{2}+x_{186}^{2}
\\ &{}+x_{190}^{2}+x_{194}^{2}-x_{197}+x_{198}^{2}-2\,x_{198}\,x_{213}+x_{199}
\\ &{}+x_{202}^{2}+x_{206}^{2}-x_{210}^{2}-x_{213}^{2}+x_{215}\,x_{221}-x_{219}
\\ &{}+x_{231}^{2}+x_{232}^{2}+x_{233}+x_{233}^{2}-4\,x_{237}-4\,x_{237}^{2}
\end{align*}
% END_Q 155
% BEGIN_Q 156
\begin{align*}
Q_{156}={}&-x_{5}\,x_{10}-x_{5}\,x_{11}-x_{5}\,x_{12}-x_{5}\,x_{13}
\\ &{}-x_{5}\,x_{14}-x_{5}\,x_{15}-x_{5}\,x_{16}-x_{5}\,x_{17}
\\ &{}-x_{5}\,x_{18}+x_{213}
\end{align*}
% END_Q 156
% BEGIN_Q 157
\begin{align*}
Q_{157}={}&-2+x_{0}-2\,x_{1}-x_{2}+2\,x_{2}\,x_{7}-x_{3}-x_{3}\,x_{10}
\\ &{}-x_{3}\,x_{11}+2\,x_{3}\,x_{21}+x_{4}-x_{5}\,x_{20}-x_{6}
\\ &{}+x_{6}^{2}+x_{11}+x_{11}^{2}+2\,x_{11}\,x_{12}+x_{12}+x_{12}^{2}
\\ &{}+x_{18}-x_{25}-3\,x_{83}-x_{84}-x_{85}-x_{144}\,x_{198}-x_{198}
\\ &{}+x_{199}-x_{214}^{2}-x_{215}^{2}-x_{216}^{2}-x_{217}^{2}-x_{218}^{2}
\\ &{}+x_{224}^{2}+x_{225}^{2}
\end{align*}
% END_Q 157
% BEGIN_Q 158
\begin{align*}
Q_{158}={}&-1+x_{2}+x_{6}-x_{6}^{2}
\end{align*}
% END_Q 158
% BEGIN_Q 159
\begin{align*}
Q_{159}={}&-x_{144}\,x_{216}-x_{216}+x_{224}
\end{align*}
% END_Q 159
% BEGIN_Q 160
\begin{align*}
Q_{160}={}&-1+x_{215}^{2}+x_{216}^{2}-x_{224}^{2}
\end{align*}
% END_Q 160
% BEGIN_Q 161
\begin{align*}
Q_{161}={}&-x_{218}-x_{218}\,x_{219}+x_{225}
\end{align*}
% END_Q 161
% BEGIN_Q 162
\begin{align*}
Q_{162}={}&-1+x_{217}^{2}+x_{218}^{2}-x_{225}^{2}
\end{align*}
% END_Q 162
% BEGIN_Q 163
\begin{align*}
Q_{163}={}&-x_{210}^{2}+x_{226}
\end{align*}
% END_Q 163
% BEGIN_Q 164
\begin{align*}
Q_{164}={}&x_{216}-2\,x_{220}\,x_{226}-x_{226}
\end{align*}
% END_Q 164
% BEGIN_Q 165
\begin{align*}
Q_{165}={}&x_{144}\,x_{215}+x_{221}
\end{align*}
% END_Q 165
% BEGIN_Q 166
\begin{align*}
Q_{166}={}&-x_{144}-x_{215}\,x_{221}+x_{219}
\end{align*}
% END_Q 166
% BEGIN_Q 167
\begin{align*}
Q_{167}={}&-1-x_{209}-x_{215}\,x_{222}+x_{217}
\end{align*}
% END_Q 167
% BEGIN_Q 168
\begin{align*}
Q_{168}={}&-x_{10}-x_{11}-4\,x_{21}-2\,x_{24}-4\,x_{26}-2\,x_{29}-2\,x_{31}
\\ &{}-4\,x_{32}-2\,x_{34}-2\,x_{36}-4\,x_{38}-2\,x_{39}-2\,x_{42}-2\,x_{46}
\\ &{}-4\,x_{67}-2\,x_{68}-x_{210}\,x_{223}+x_{218}
\end{align*}
% END_Q 168
% BEGIN_Q 169
\begin{align*}
Q_{169}={}&-1+2\,x_{0}^{2}+2\,x_{0}\,x_{4}-2\,x_{0}\,x_{149}-2\,x_{0}\,x_{158}+x_{1}^{2}
\\ &{}-2\,x_{1}\,x_{198}+x_{2}^{2}-2\,x_{2}\,x_{186}+x_{4}^{2}-2\,x_{4}\,x_{158}
\\ &{}+x_{10}^{2}-2\,x_{10}\,x_{164}-2\,x_{10}\,x_{210}+x_{11}^{2}-2\,x_{11}\,x_{170}
\\ &{}-2\,x_{11}\,x_{210}+x_{12}^{2}-2\,x_{12}\,x_{174}+x_{13}^{2}-2\,x_{13}\,x_{178}
\\ &{}+x_{14}^{2}-2\,x_{14}\,x_{182}+x_{15}^{2}-2\,x_{15}\,x_{190}+x_{16}^{2}
\\ &{}-2\,x_{16}\,x_{194}+x_{17}^{2}-2\,x_{17}\,x_{202}+x_{18}^{2}-2\,x_{18}\,x_{206}
\\ &{}-8\,x_{21}\,x_{210}-4\,x_{24}\,x_{210}-8\,x_{26}\,x_{210}-4\,x_{29}\,x_{210}
\\ &{}-4\,x_{31}\,x_{210}-8\,x_{32}\,x_{210}-4\,x_{34}\,x_{210}-4\,x_{36}\,x_{210}
\\ &{}-8\,x_{38}\,x_{210}-4\,x_{39}\,x_{210}-4\,x_{42}\,x_{210}-4\,x_{46}\,x_{210}
\\ &{}-8\,x_{67}\,x_{210}-4\,x_{68}\,x_{210}+x_{144}-x_{144}\,x_{164}-x_{144}\,x_{227}
\\ &{}+x_{148}^{2}+x_{157}^{2}+x_{163}^{2}+x_{164}+x_{165}+x_{169}^{2}
\\ &{}+x_{173}^{2}+x_{177}^{2}+x_{181}^{2}+x_{185}^{2}+x_{189}^{2}+x_{193}^{2}
\\ &{}+x_{197}^{2}+x_{201}^{2}+x_{205}^{2}+x_{209}^{2}+x_{210}^{2}+x_{226}
\\ &{}-x_{231}^{2}-x_{232}^{2}-x_{233}-x_{233}^{2}
\end{align*}
% END_Q 169
% BEGIN_Q 170
\begin{align*}
Q_{170}={}&-x_{105}\,x_{228}+x_{130}
\end{align*}
% END_Q 170
% BEGIN_Q 171
\begin{align*}
Q_{171}={}&-x_{110}\,x_{229}+x_{130}
\end{align*}
% END_Q 171
% BEGIN_Q 172
\begin{align*}
Q_{172}={}&-x_{123}\,x_{230}+x_{130}
\end{align*}
% END_Q 172
% BEGIN_Q 173
\begin{align*}
Q_{173}={}&2+x_{1}-x_{2}\,x_{7}-x_{3}\,x_{10}-x_{3}\,x_{11}-x_{4}
\\ &{}+x_{5}\,x_{20}+x_{10}+x_{10}^{2}+2\,x_{10}\,x_{11}-8\,x_{10}\,x_{21}
\\ &{}-4\,x_{10}\,x_{24}-8\,x_{10}\,x_{26}-4\,x_{10}\,x_{29}-4\,x_{10}\,x_{31}
\\ &{}-8\,x_{10}\,x_{32}-4\,x_{10}\,x_{34}-4\,x_{10}\,x_{36}-8\,x_{10}\,x_{38}
\\ &{}-4\,x_{10}\,x_{39}-4\,x_{10}\,x_{42}-4\,x_{10}\,x_{46}-8\,x_{10}\,x_{67}
\\ &{}-4\,x_{10}\,x_{68}+2\,x_{11}+3\,x_{11}^{2}-2\,x_{11}\,x_{12}-8\,x_{11}\,x_{21}
\\ &{}-4\,x_{11}\,x_{24}-8\,x_{11}\,x_{26}-4\,x_{11}\,x_{29}-4\,x_{11}\,x_{31}
\\ &{}-8\,x_{11}\,x_{32}-4\,x_{11}\,x_{34}-4\,x_{11}\,x_{36}+4\,x_{11}\,x_{38}
\\ &{}-4\,x_{11}\,x_{39}-4\,x_{11}\,x_{42}-4\,x_{11}\,x_{46}-8\,x_{11}\,x_{67}
\\ &{}-4\,x_{11}\,x_{68}-x_{12}-x_{12}^{2}-x_{18}+8\,x_{21}-16\,x_{21}\,x_{24}
\\ &{}-32\,x_{21}\,x_{26}-16\,x_{21}\,x_{29}-16\,x_{21}\,x_{31}-32\,x_{21}\,x_{32}
\\ &{}-16\,x_{21}\,x_{34}-16\,x_{21}\,x_{36}-32\,x_{21}\,x_{38}-16\,x_{21}\,x_{39}
\\ &{}-16\,x_{21}\,x_{42}-16\,x_{21}\,x_{46}-32\,x_{21}\,x_{67}-16\,x_{21}\,x_{68}
\\ &{}-2\,x_{24}-4\,x_{24}^{2}-16\,x_{24}\,x_{26}-8\,x_{24}\,x_{29}-8\,x_{24}\,x_{31}
\\ &{}-16\,x_{24}\,x_{32}-8\,x_{24}\,x_{34}-8\,x_{24}\,x_{36}-16\,x_{24}\,x_{38}
\\ &{}-8\,x_{24}\,x_{39}-8\,x_{24}\,x_{42}-8\,x_{24}\,x_{46}-16\,x_{24}\,x_{67}
\\ &{}-8\,x_{24}\,x_{68}+x_{25}-4\,x_{26}-12\,x_{26}^{2}+8\,x_{26}\,x_{28}
\\ &{}-16\,x_{26}\,x_{29}-16\,x_{26}\,x_{31}-32\,x_{26}\,x_{32}-16\,x_{26}\,x_{34}
\\ &{}-16\,x_{26}\,x_{36}-32\,x_{26}\,x_{38}-16\,x_{26}\,x_{39}-16\,x_{26}\,x_{42}
\\ &{}-16\,x_{26}\,x_{46}-32\,x_{26}\,x_{67}-16\,x_{26}\,x_{68}+4\,x_{28}^{2}-3\,x_{29}
\\ &{}-4\,x_{29}^{2}-8\,x_{29}\,x_{31}-16\,x_{29}\,x_{32}-8\,x_{29}\,x_{34}-8\,x_{29}\,x_{36}
\\ &{}-16\,x_{29}\,x_{38}-8\,x_{29}\,x_{39}-8\,x_{29}\,x_{42}-8\,x_{29}\,x_{46}
\\ &{}-16\,x_{29}\,x_{67}-8\,x_{29}\,x_{68}-2\,x_{31}-4\,x_{31}^{2}-16\,x_{31}\,x_{32}
\\ &{}-8\,x_{31}\,x_{34}-8\,x_{31}\,x_{36}-16\,x_{31}\,x_{38}-8\,x_{31}\,x_{39}
\\ &{}-8\,x_{31}\,x_{42}-8\,x_{31}\,x_{46}-16\,x_{31}\,x_{67}-8\,x_{31}\,x_{68}-4\,x_{32}
\\ &{}-12\,x_{32}^{2}+8\,x_{32}\,x_{33}-16\,x_{32}\,x_{34}-16\,x_{32}\,x_{36}
\\ &{}-32\,x_{32}\,x_{38}-16\,x_{32}\,x_{39}-16\,x_{32}\,x_{42}-16\,x_{32}\,x_{46}
\\ &{}-32\,x_{32}\,x_{67}-16\,x_{32}\,x_{68}+4\,x_{33}^{2}-3\,x_{34}-4\,x_{34}^{2}
\\ &{}-8\,x_{34}\,x_{36}-16\,x_{34}\,x_{38}-8\,x_{34}\,x_{39}-8\,x_{34}\,x_{42}
\\ &{}-8\,x_{34}\,x_{46}-16\,x_{34}\,x_{67}-8\,x_{34}\,x_{68}-2\,x_{36}-4\,x_{36}^{2}
\\ &{}-16\,x_{36}\,x_{38}-8\,x_{36}\,x_{39}-8\,x_{36}\,x_{42}-8\,x_{36}\,x_{46}
\\ &{}-16\,x_{36}\,x_{67}-8\,x_{36}\,x_{68}-4\,x_{38}-4\,x_{38}^{2}-16\,x_{38}\,x_{39}
\\ &{}-16\,x_{38}\,x_{42}-16\,x_{38}\,x_{46}-32\,x_{38}\,x_{67}-16\,x_{38}\,x_{68}
\\ &{}-5\,x_{39}-4\,x_{39}^{2}-8\,x_{39}\,x_{42}-8\,x_{39}\,x_{46}-16\,x_{39}\,x_{67}
\\ &{}-8\,x_{39}\,x_{68}-x_{40}-2\,x_{42}-4\,x_{42}^{2}-8\,x_{42}\,x_{46}
\\ &{}-16\,x_{42}\,x_{67}-8\,x_{42}\,x_{68}-2\,x_{46}-4\,x_{46}^{2}-16\,x_{46}\,x_{67}
\\ &{}-8\,x_{46}\,x_{68}+4\,x_{67}-8\,x_{67}^{2}-16\,x_{67}\,x_{68}-4\,x_{68}-4\,x_{68}^{2}
\\ &{}+x_{83}+x_{84}-x_{149}^{2}-x_{158}^{2}+2\,x_{164}-x_{164}^{2}-x_{170}^{2}
\\ &{}-x_{174}^{2}-x_{178}^{2}-x_{182}^{2}-x_{186}^{2}-x_{190}^{2}-x_{194}^{2}
\\ &{}-x_{198}^{2}-x_{202}^{2}-x_{206}^{2}+x_{210}^{2}-x_{210}\,x_{223}+x_{218}
\\ &{}-x_{231}^{2}-x_{232}^{2}-x_{233}-x_{233}^{2}
\end{align*}
% END_Q 173
% BEGIN_Q 174
\begin{align*}
Q_{174}={}&-1-4\,x_{107}^{2}-4\,x_{112}^{2}-4\,x_{115}^{2}-4\,x_{118}^{2}-4\,x_{121}^{2}-4\,x_{125}^{2}
\\ &{}-4\,x_{128}^{2}+x_{130}^{2}-4\,x_{132}^{2}-4\,x_{136}^{2}-4\,x_{139}^{2}-4\,x_{142}^{2}
\\ &{}-x_{234}^{2}-x_{235}^{2}-x_{236}-x_{236}^{2}-4\,x_{239}^{2}
\end{align*}
% END_Q 174
% BEGIN_Q 175
\begin{align*}
Q_{175}={}&-x_{5}\,x_{238}+x_{130}
\end{align*}
% END_Q 175
% BEGIN_Q 176
\begin{align*}
Q_{176}={}&x_{3}\,x_{130}-x_{6}\,x_{238}-x_{239}
\end{align*}
% END_Q 176
\endgroup
\section{Coordinate names and artifact interpretation}
The semantic names below describe the restored graph coordinates; the equation itself uses only the independent \(x_i\). Recorded translation shifts must be applied before identifying an actual source quantity with a final coordinate. The names are annotations, not new equations.
\begingroup\footnotesize\renewcommand{\arraystretch}{0.96}
\begin{longtable}{r >{\raggedright\arraybackslash}p{0.38\textwidth} r >{\raggedright\arraybackslash}p{0.38\textwidth}}
\toprule Index & Source name & Index & Source name\\\midrule\endhead
0 & \path{t0} & 120 & \path{range3.quotient}\\
1 & \path{x} & 121 & \path{normalized.range3.remainder}\\
2 & \path{y} & 122 & \path{pow10.output}\\
3 & \path{fact0.output} & 123 & \path{pow11.output}\\
4 & \path{mul.1} & 124 & \path{range4.quotient}\\
5 & \path{choose0.output} & 125 & \path{normalized.range4.remainder}\\
6 & \path{pow0.output} & 126 & \path{pow12.output}\\
7 & \path{endpoint.u0} & 127 & \path{range5.quotient}\\
8 & \path{endpoint.ut} & 128 & \path{normalized.range5.remainder}\\
9 & \path{mul.11} & 129 & \path{pow13.output}\\
10 & \path{crt.encoded.a} & 130 & \path{pow14.output}\\
11 & \path{crt.encoded.b} & 131 & \path{range6.quotient}\\
12 & \path{crt.encoded.c} & 132 & \path{normalized.range6.remainder}\\
13 & \path{crt.encoded.d} & 133 & \path{mul.430}\\
14 & \path{crt.encoded.S} & 134 & \path{pow15.output}\\
15 & \path{crt.encoded.Sp} & 135 & \path{range7.quotient}\\
16 & \path{crt.encoded.u} & 136 & \path{normalized.range7.remainder}\\
17 & \path{crt.encoded.vq} & 137 & \path{pow16.output}\\
18 & \path{crt.encoded_index_minus_one} & 138 & \path{range8.quotient}\\
19 & \path{pow1.output} & 139 & \path{normalized.range8.remainder}\\
20 & \path{crt.index_quotient} & 140 & \path{pow17.output}\\
21 & \path{mul.40} & 141 & \path{range9.quotient}\\
22 & \path{add.50} & 142 & \path{normalized.range9.remainder}\\
23 & \path{mul.53} & 143 & \path{shared.W}\\
24 & \path{add.61} & 144 & \path{shared.A}\\
25 & \path{mul.64} & 145 & \path{shared.Zminus1}\\
26 & \path{add.71} & 146 & \path{mul.450}\\
27 & \path{mul.75} & 147 & \path{mul.451}\\
28 & \path{add.83} & 148 & \path{index0.X}\\
29 & \path{mul.86} & 149 & \path{index0.Y}\\
30 & \path{mul.96} & 150 & \path{mul.460}\\
31 & \path{add.101} & 151 & \path{base0.T}\\
32 & \path{add.112} & 152 & \path{pow0.quotient}\\
33 & \path{add.119} & 153 & \path{base1.T}\\
34 & \path{mul.123} & 154 & \path{pow1.quotient}\\
35 & \path{mul.137} & 155 & \path{base2.T}\\
36 & \path{add.142} & 156 & \path{pow2.quotient}\\
37 & \path{mul.148} & 157 & \path{index2.X}\\
38 & \path{add.156} & 158 & \path{index2.Y}\\
39 & \path{mul.159} & 159 & \path{mul.503}\\
40 & \path{add.172} & 160 & \path{base3.T}\\
41 & \path{add.177} & 161 & \path{pow3.quotient}\\
42 & \path{add.181} & 162 & \path{pow4.quotient}\\
43 & \path{mul.184} & 163 & \path{index3.X}\\
44 & \path{mul.188} & 164 & \path{index3.Y}\\
45 & \path{mul.194} & 165 & \path{mul.529}\\
46 & \path{add.199} & 166 & \path{pow5.quotient}\\
47 & \path{add.205} & 167 & \path{index_addition.D}\\
48 & \path{mul.210} & 168 & \path{pow6.quotient}\\
49 & \path{mul.214} & 169 & \path{index5.X}\\
50 & \path{add.221} & 170 & \path{index5.Y}\\
51 & \path{mul.225} & 171 & \path{mul.554}\\
52 & \path{mul.229} & 172 & \path{pow7.quotient}\\
53 & \path{add.236} & 173 & \path{index6.X}\\
54 & \path{mul.240} & 174 & \path{index6.Y}\\
55 & \path{add.247} & 175 & \path{mul.567}\\
56 & \path{mul.251} & 176 & \path{pow8.quotient}\\
57 & \path{mul.255} & 177 & \path{index7.X}\\
58 & \path{add.262} & 178 & \path{index7.Y}\\
59 & \path{mul.266} & 179 & \path{mul.580}\\
60 & \path{mul.270} & 180 & \path{pow9.quotient}\\
61 & \path{add.277} & 181 & \path{index8.X}\\
62 & \path{mul.281} & 182 & \path{index8.Y}\\
63 & \path{mul.285} & 183 & \path{mul.593}\\
64 & \path{add.292} & 184 & \path{pow10.quotient}\\
65 & \path{add.299} & 185 & \path{index9.X}\\
66 & \path{mul.303} & 186 & \path{index9.Y}\\
67 & \path{add.308} & 187 & \path{mul.606}\\
68 & \path{mul.311} & 188 & \path{pow11.quotient}\\
69 & \path{mul.314} & 189 & \path{index10.X}\\
70 & \path{sqd.1} & 190 & \path{index10.Y}\\
71 & \path{sqd.2} & 191 & \path{mul.619}\\
72 & \path{sqd.3} & 192 & \path{pow12.quotient}\\
73 & \path{sqd.4} & 193 & \path{index11.X}\\
74 & \path{sqd.5} & 194 & \path{index11.Y}\\
75 & \path{sqd.6} & 195 & \path{mul.632}\\
76 & \path{sqd.7} & 196 & \path{pow13.quotient}\\
77 & \path{sqd.8} & 197 & \path{index12.X}\\
78 & \path{sqd.9} & 198 & \path{index12.Y}\\
79 & \path{mul.331} & 199 & \path{mul.645}\\
80 & \path{sqd.11} & 200 & \path{pow14.quotient}\\
81 & \path{mul.334} & 201 & \path{index13.X}\\
82 & \path{sqd.12} & 202 & \path{index13.Y}\\
83 & \path{mul.337} & 203 & \path{mul.658}\\
84 & \path{mul.340} & 204 & \path{pow15.quotient}\\
85 & \path{mul.342} & 205 & \path{index14.X}\\
86 & \path{mul.344} & 206 & \path{index14.Y}\\
87 & \path{sqd.13} & 207 & \path{mul.671}\\
88 & \path{sqd.14} & 208 & \path{pow16.quotient}\\
89 & \path{mul.350} & 209 & \path{index15.X}\\
90 & \path{mul.351} & 210 & \path{index15.Y}\\
91 & \path{mul.352} & 211 & \path{mul.684}\\
92 & \path{mul.353} & 212 & \path{pow17.quotient}\\
93 & \path{mul.354} & 213 & \path{self_radix.BH}\\
94 & \path{mul.355} & 214 & \path{cofinal.horizon_root}\\
95 & \path{mul.356} & 215 & \path{cap.H}\\
96 & \path{mul.357} & 216 & \path{cap.J}\\
97 & \path{mul.358} & 217 & \path{cap.S}\\
98 & \path{mul.359} & 218 & \path{cap.R}\\
99 & \path{mul.360} & 219 & \path{cap.B}\\
100 & \path{mul.362} & 220 & \path{cap.e}\\
101 & \path{mul.364} & 221 & \path{cap.f}\\
102 & \path{mul.367} & 222 & \path{cap.g}\\
103 & \path{crt.nonnegative_quotient} & 223 & \path{cap.h}\\
104 & \path{pow3.output} & 224 & \path{cap.common_Ay}\\
105 & \path{pow4.output} & 225 & \path{cap.BR}\\
106 & \path{choose0.high} & 226 & \path{cap.Y2}\\
107 & \path{normalized.choose0.remainder} & 227 & \path{cap.pooled_quotient}\\
108 & \path{mul.396} & 228 & \path{remainder_budget.ratio0}\\
109 & \path{pow5.output} & 229 & \path{remainder_budget.ratio1}\\
110 & \path{pow6.output} & 230 & \path{remainder_budget.ratio2}\\
111 & \path{range0.quotient} & 231 & \path{new_guard3.a}\\
112 & \path{normalized.range0.remainder} & 232 & \path{new_guard3.b}\\
113 & \path{pow7.output} & 233 & \path{new_guard3.c}\\
114 & \path{range1.quotient} & 234 & \path{new_guard4.a}\\
115 & \path{normalized.range1.remainder} & 235 & \path{new_guard4.b}\\
116 & \path{pow8.output} & 236 & \path{new_guard4.c}\\
117 & \path{range2.quotient} & 237 & \path{radix.odd_square_root_half}\\
118 & \path{normalized.range2.remainder} & 238 & \path{bootstrap.normalizing_ratio}\\
119 & \path{pow9.output} & 239 & \path{bootstrap.normalized_ceiling_error}\\
\bottomrule\end{longtable}\endgroup
\end{document}
